%% titre: Cours d'algèbre 1
%% matiere: Algèbre 1
%% niveau: Licence 1
%% annee: 2026-2027
%% description: Cours d'algèbre 1 Version 1.
%% ordre: 1
\documentclass[a4paper,12pt,openany]{book}

\usepackage[utf8]{inputenc}
\usepackage[T1]{fontenc}
\usepackage[french,provide=*]{babel}

\providecommand{\og}{\guillemotleft~}
\providecommand{\fg}{\ifhmode~\fi\guillemotright}
\providecommand{\up}[1]{\textsuperscript{#1}}
\usepackage{geometry}
\usepackage{amsmath,amssymb,amsthm}
\usepackage{cancel}
\usepackage{enumitem}
\usepackage{array}
\usepackage{booktabs}
\usepackage{longtable}
\usepackage{tabularx}
\usepackage{tcolorbox}
\usepackage{titlesec}
\usepackage{fancyhdr}
%\usepackage{hyperref}
\usepackage{xcolor}
%\usepackage{tikz}
\usepackage{enumitem}
% ============================================================

\newtheorem{application}{Application}[section]
\newtheorem{resume}{Resumé}[section]
\begin{document}
\tableofcontents
%\newpage

% ============================================================
% CHAPITRE 1
%
%\thispagestyle{empty}
%\vspace*{\fill}
%\begin{center}
%	{\Huge Chapitre 1}
%\end{center}
%\vspace*{\fill}
%\newpage
% ============================================================

\chapter{Notions de Logique}

\section{Introduction}


Ce premier chapitre a pour objectifs :
\begin{itemize}
  \item d'apprendre à reconnaître et à manipuler des \textbf{propositions} (des
  énoncés vrais ou faux) ;
  \item de construire des propositions composées à l'aide des
  \textbf{connecteurs logiques} (négation, conjonction, disjonction,
  implication, équivalence) ;
  \item d'introduire les \textbf{prédicats} et les \textbf{quantificateurs}
  ($\forall$, $\exists$), indispensables dès qu'on manipule des variables ;
  \item de présenter les grands \textbf{types de raisonnement} utilisés tout au
  long des études mathématiques : raisonnement déductif, par contraposée, par
  l'absurde, par disjonction de cas, et par récurrence.
\end{itemize}

Ce chapitre est le socle de tout le reste du cours : les chapitres suivants
(ensembles, applications, structures algébriques) utiliseront en permanence le
vocabulaire et les techniques introduits ici. On ne peut pas comprendre une
démonstration de théorème si l'on ne maîtrise pas le sens exact d'une
implication ou d'un \og pour tout \fg{}.

\begin{application}
La logique n'est pas qu'un outil de mathématicien : elle est à la base de la
\textbf{conception des circuits électroniques} (portes logiques ET, OU, NON),
de la \textbf{programmation informatique} (conditions \texttt{if}, boucles,
expressions booléennes), des \textbf{bases de données} (requêtes SQL avec
\texttt{AND}, \texttt{OR}, \texttt{NOT}) et de l'\textbf{intelligence
artificielle} (systèmes experts, logique des prédicats en IA symbolique).
\end{application}

\section{Prérequis}

Ce chapitre ne suppose \emph{aucune} connaissance mathématique préalable au-delà
du niveau du baccalauréat scientifique : savoir manipuler des nombres entiers,
des égalités et des inégalités simples. Tout le vocabulaire logique est
introduit progressivement et défini avant d'être utilisé.

\section{Propositions et connecteurs logiques}

\subsection{Notion de proposition}


\begin{definition}
On appelle \textbf{proposition} tout énoncé qui est soit vrai, soit faux (et
jamais les deux à la fois).
\end{definition}

Les propositions sont en général notées par des lettres majuscules : $P$, $Q$,
$R$, $P_0$, etc. Lorsqu'une proposition $P$ est vraie, on dit \og on a $P$ \fg.

\begin{exemple}
\begin{enumerate}
  \item \og $3+5=8$ \fg{} est une proposition : elle est vraie.
  \item \og $2\times 3 = 5$ \fg{} est une proposition : elle est fausse.
  \item \og Comment allez-vous ? \fg{} n'est pas une proposition : ce n'est ni
  vrai ni faux, c'est une question.
  \item \og L'entier $a$ divise $2$ \fg{} n'est pas une proposition \emph{telle
  quelle}, car sa vérité dépend de la valeur de $a$ (c'est un \emph{prédicat},
  notion que nous verrons en section~1.4). Ce n'est une proposition que
  lorsqu'on remplace $a$ par un entier précis : \og $4$ divise $2$ \fg{} est
  une proposition (fausse), \og $1$ divise $2$ \fg{} en est une autre (vraie).
\end{enumerate}
\end{exemple}


À partir d'une ou plusieurs propositions, on construit de nouvelles
propositions, dites \textbf{composées}, à l'aide de \textbf{connecteurs
logiques}. Nous en étudions cinq : la négation, la conjonction, la disjonction,
l'implication et l'équivalence.

\subsection{La négation}


\begin{definition}
La \textbf{négation} d'une proposition $P$ est la proposition, notée $\overline{P}$
(ou $\lnot P$), qui est vraie si $P$ est fausse, et fausse si $P$ est vraie.
$\overline{P}$ se lit \og non $P$ \fg.
\end{definition}

\begin{exemple}
Soit $P$ : \og Je suis étudiant \fg. Alors $\overline{P}$ : \og Je ne suis pas
étudiant \fg.

Soit $Q$ : \og ce tableau est rouge \fg. Alors $\overline{Q}$ : \og ce tableau
n'est pas rouge \fg.
\end{exemple}


On peut résumer cette définition par un tableau, appelé \textbf{table de
vérité}, où l'on convient de noter $1$ la valeur \og vrai \fg{} et $0$ la
valeur \og faux \fg :

\begin{center}
\begin{tabular}{|c|c|}
\hline
$P$ & $\overline{P}$ \\
\hline
1 & 0 \\
0 & 1 \\
\hline
\end{tabular}
\end{center}

\subsection{La conjonction}

\begin{definition}
La \textbf{conjonction} de deux propositions $P$ et $Q$ est la proposition,
notée $P\wedge Q$, qui est vraie lorsque $P$ et $Q$ sont \emph{simultanément}
vraies, et fausse dans tous les autres cas. $P\wedge Q$ se lit \og $P$ et
$Q$ \fg.
\end{definition}

\begin{center}
\begin{tabular}{|c|c|c|}
\hline
$P$ & $Q$ & $P\wedge Q$ \\
\hline
1 & 1 & 1 \\
1 & 0 & 0 \\
0 & 1 & 0 \\
0 & 0 & 0 \\
\hline
\end{tabular}
\end{center}

\begin{exemple}
Soit $P$ : \og $6$ est pair \fg{} (vraie) et $Q$ : \og $6$ est un multiple de
$5$ \fg{} (fausse). Alors $P\wedge Q$ : \og $6$ est pair \emph{et} $6$ est un
multiple de $5$ \fg{} est \textbf{fausse}, car $Q$ est fausse.
\end{exemple}

\subsection{La disjonction}

\begin{definition}
La \textbf{disjonction} de deux propositions $P$ et $Q$ est la proposition,
notée $P\vee Q$, qui est fausse lorsque $P$ et $Q$ sont \emph{simultanément}
fausses, et vraie dans tous les autres cas. $P\vee Q$ se lit \og $P$ ou
$Q$ \fg.
\end{definition}

\begin{center}
\begin{tabular}{|c|c|c|}
\hline
$P$ & $Q$ & $P\vee Q$ \\
\hline
1 & 1 & 1 \\
1 & 0 & 1 \\
0 & 1 & 1 \\
0 & 0 & 0 \\
\hline
\end{tabular}
\end{center}


\begin{exemple}
Soit $P$ : \og $7$ est premier \fg{} (vraie) et $Q$ : \og $7$ est pair \fg{}
(fausse). Alors $P\vee Q$ est \textbf{vraie}, car il suffit que l'une des deux
propositions soit vraie.
\end{exemple}

\subsection{L'implication}


\begin{definition}
L'\textbf{implication} de deux propositions $P$ et $Q$ (dans cet ordre) est la
proposition, notée $P\Rightarrow Q$, qui est \textbf{fausse} uniquement lorsque
$P$ est vraie et $Q$ est fausse, et \textbf{vraie} dans tous les autres cas.
$P\Rightarrow Q$ se lit \og $P$ implique $Q$ \fg{} ou \og si $P$ alors $Q$ \fg.
\end{definition}

\begin{center}
\begin{tabular}{|c|c|c|}
\hline
$P$ & $Q$ & $P\Rightarrow Q$ \\
\hline
1 & 1 & 1 \\
1 & 0 & 0 \\
0 & 1 & 1 \\
0 & 0 & 1 \\
\hline
\end{tabular}
\end{center}

\begin{remarque}[Pourquoi $P\Rightarrow Q$ est-elle vraie quand $P$ est fausse ?]
C'est le point le plus délicat pour un débutant. Intuitivement, une implication
\og Si $P$ alors $Q$ \fg{} promet seulement quelque chose \emph{au cas où} $P$
serait vraie ; elle ne dit rien sur ce qui se passe si $P$ ne se réalise pas.
Une promesse non testée ne peut pas être prise en défaut : on considère donc
qu'elle est vraie par défaut. Par exemple, la phrase \og S'il pleut, alors je
prends un parapluie \fg{} n'est jamais mise en défaut les jours où il ne pleut
pas, quel que soit mon comportement.
\end{remarque}

\begin{exemple}
Soit $P$ : \og $3=7-4$ \fg{} (vraie) et $Q$ : \og $7=4-3$ \fg{} (fausse).
\begin{enumerate}
  \item $P\Rightarrow Q$ : \og $3=7-4$ implique $7=4-3$ \fg{} est
  \textbf{fausse} (car $P$ vraie et $Q$ fausse).
  \item $Q\Rightarrow P$ : \og $7=4-3$ implique $3=7-4$ \fg{} est
  \textbf{vraie} (car $Q$ fausse : troisième ou quatrième ligne de la table).
\end{enumerate}
\end{exemple}

\begin{exemple}
La proposition \og Si $2=3$, alors $1=0$ \fg{} est \textbf{vraie}, bien que ni
\og $2=3$ \fg{} ni \og $1=0$ \fg{} ne soient vraies individuellement. En effet,
l'hypothèse \og $2=3$ \fg{} est fausse, donc l'implication est automatiquement
vraie, d'après la définition.
\end{exemple}


\subsection{L'équivalence}

\begin{definition}
L'\textbf{équivalence} de deux propositions $P$ et $Q$ est la proposition,
notée $P\Leftrightarrow Q$, qui est vraie lorsque $P$ et $Q$ ont la
\emph{même} valeur de vérité (toutes deux vraies, ou toutes deux fausses), et
fausse sinon. $P\Leftrightarrow Q$ se lit \og $P$ équivaut à $Q$ \fg{} ou
\og $P$ si et seulement si $Q$ \fg.
\end{definition}

\begin{center}
\begin{tabular}{|c|c|c|}
\hline
$P$ & $Q$ & $P\Leftrightarrow Q$ \\
\hline
1 & 1 & 1 \\
1 & 0 & 0 \\
0 & 1 & 0 \\
0 & 0 & 1 \\
\hline
\end{tabular}
\end{center}

\begin{exemple}
Soit $P$ : \og $x=2$ \fg{} et $Q$ : \og $x^2=4$ et $x>0$ \fg. Pour tout réel
$x$, $P\Leftrightarrow Q$ est vraie : dès que $x=2$, on a bien $x^2=4$ et
$x>0$ ; et réciproquement, si $x^2=4$ et $x>0$, alors nécessairement $x=2$.
\end{exemple}

\subsection{Tableau récapitulatif}

En regroupant les cinq connecteurs, on obtient le tableau de vérité suivant,
qui synthétise tout ce que nous avons vu :

\begin{center}
\begin{tabular}{|c|c|c|c|c|c|c|}
\hline
$P$ & $Q$ & $\overline{P}$ & $P\wedge Q$ & $P\vee Q$ & $P\Rightarrow Q$ & $P\Leftrightarrow Q$\\
\hline
1 & 1 & 0 & 1 & 1 & 1 & 1 \\
1 & 0 & 0 & 0 & 1 & 0 & 0 \\
0 & 1 & 1 & 0 & 1 & 1 & 0 \\
0 & 0 & 1 & 0 & 0 & 1 & 1 \\
\hline
\end{tabular}
\end{center}

\section{Calcul des prédicats}

\subsection{Notion de prédicat}


\begin{definition}
On appelle \textbf{prédicat} tout énoncé contenant une ou plusieurs variables,
qui devient une proposition dès que l'on remplace les variables par des objets
concrets (précis). Les prédicats sont souvent notés $P(x)$, $Q(x,y)$, $R(z)$,
etc.
\end{definition}

\begin{exemple}
\begin{enumerate}
  \item $P(x)$ : \og l'entier $x$ divise $2$ \fg{} est un prédicat. On a par
  exemple $P(1)$ vraie et $P(3)$ fausse.
  \item $Q(x)$ : \og le réel $x$ est supérieur à $0$ \fg{} est un prédicat.
  \item $R(n,m)$ : \og la différence des entiers $n$ et $m$ est un multiple de
  $3$ \fg{} est un prédicat à deux variables.
\end{enumerate}
\end{exemple}

On peut combiner des prédicats avec les connecteurs logiques déjà vus, exactement
comme pour des propositions : $P(x)\wedge Q(x)$, $P(x)\Rightarrow Q(x)$, etc.
sont encore des prédicats.

\subsection{Les quantificateurs}


\begin{definition}
Soit $P(x)$ un prédicat et $E$ un ensemble non vide.
\begin{enumerate}
  \item L'expression \og pour tout élément $x$ de $E$, $P(x)$ est vraie \fg{}
  s'écrit en abrégé : $\forall x \in E,\ P(x)$. Le symbole $\forall$ s'appelle
  le \textbf{quantificateur universel}.
  \item L'expression \og il existe au moins un élément $x$ de $E$ tel que
  $P(x)$ soit vraie \fg{} s'écrit en abrégé : $\exists x \in E,\ P(x)$. Le
  symbole $\exists$ s'appelle le \textbf{quantificateur existentiel}.
\end{enumerate}
\end{definition}

\begin{exemple}
\begin{enumerate}
  \item \og $\forall x \in \mathbb{Z},\ x = x^3$ \fg{} signifie \og pour tout
  entier $x$, on a $x=x^3$ \fg{}. Cette proposition est \textbf{fausse} : par
  exemple $2 \neq 2^3=8$.
  \item \og $\exists x \in \mathbb{N},\ x-4>0$ \fg{} signifie \og il existe un
  entier naturel $x$ tel que $x-4>0$ \fg. Cette proposition est
  \textbf{vraie} : $x=5$ convient, car $5-4=1>0$.
  \item \og $\exists x \in \mathbb{R},\ x^2<0$ \fg{} est \textbf{fausse} :
  aucun carré de réel n'est strictement négatif.
  \item \og $\forall a \in \mathbb{N},\ 1 \text{ divise } a$ \fg{} est
  \textbf{vraie} : tout entier naturel est divisible par $1$.
\end{enumerate}
\end{exemple}

\begin{remarque}
\og $\forall x \in E,\ P(x)$ \fg{} se lit aussi \og quel que soit $x$ de
$E$, $P(x)$ \fg.
\end{remarque}

\paragraph{Prédicats à deux variables et deux quantificateurs.}
Lorsqu'un prédicat $P(x,y)$ dépend de deux variables $x\in E$ et $y\in F$, on
peut le quantifier successivement. \textbf{L'ordre des quantificateurs a une
importance capitale} lorsqu'ils ne sont pas de même nature.

\begin{definition}
Soient $E$ et $F$ deux ensembles et $P(x,y)$ un prédicat à deux variables.
\begin{itemize}
  \item $\forall x \in E,\ \forall y \in F,\ P(x,y)$ signifie que pour
  \emph{tout} $x$ de $E$ et \emph{tout} $y$ de $F$, $P(x,y)$ est vraie.
  \item $\exists x \in E,\ \exists y \in F,\ P(x,y)$ signifie qu'il existe
  \emph{au moins un} $x$ de $E$ et \emph{au moins un} $y$ de $F$ tels que
  $P(x,y)$ soit vraie.
  \item $\forall x \in E,\ \exists y \in F,\ P(x,y)$ signifie que pour
  \emph{chaque} $x$ de $E$, il existe (au moins) un $y$ de $F$, qui peut
  \textbf{dépendre de $x$}, tel que $P(x,y)$ soit vraie.
  \item $\exists y \in F,\ \forall x \in E,\ P(x,y)$ signifie qu'il existe un
  \emph{même} $y$ de $F$, \textbf{indépendant de $x$}, qui convient pour
  \emph{tous} les $x$ de $E$.
\end{itemize}
\end{definition}


\begin{exemple}
\begin{enumerate}
  \item \og $\forall x \in \mathbb{N},\ \exists y \in \mathbb{R},\ x \leq y$
  \fg{} est \textbf{vraie} : pour chaque entier $x$, on peut choisir $y=x$
  (le choix de $y$ dépend de $x$, ce qui est autorisé).
  \item \og $\exists x \in \mathbb{N},\ \forall y \in \mathbb{R},\ x \leq y$
  \fg{} est \textbf{fausse} : il n'existe pas d'entier $x$ qui soit inférieur
  ou égal à \emph{tous} les réels $y$ (en particulier aux réels négatifs).
\end{enumerate}
\end{exemple}

\subsection{Négation des propositions quantifiées}

\begin{proposition}
Soit $P(x)$ un prédicat sur un ensemble $E$. On a les règles de négation
suivantes :
\[
\overline{\forall x \in E,\ P(x)} \iff \exists x \in E,\ \overline{P(x)}
\qquad\text{et}\qquad
\overline{\exists x \in E,\ P(x)} \iff \forall x \in E,\ \overline{P(x)}.
\]
Autrement dit, pour nier une phrase quantifiée, on \textbf{échange
$\forall$ et $\exists$}, et on nie le prédicat final.
\end{proposition}

\begin{remarque}[Intuition]
Dire que \og tous les éléments vérifient $P$ \fg{} est faux revient exactement
à dire qu'\og il en existe au moins un qui ne vérifie pas $P$ \fg. De même, dire
qu'\og il en existe un qui vérifie $P$ \fg{} est faux revient à dire
qu'\og aucun ne le vérifie \fg, c'est-à-dire que \og tous ne le vérifient pas \fg.
\end{remarque}

\begin{exemple}
\begin{enumerate}
  \item La négation de \og $\forall x \in \mathbb{Z},\ x=x^3$ \fg{} est
  \og $\exists x \in \mathbb{Z},\ x \neq x^3$ \fg{} (proposition vraie, car
  $x=2$ convient).
  \item La négation de \og $\exists x \in \mathbb{N},\ x-4>0$ \fg{} est
  \og $\forall x \in \mathbb{N},\ x-4 \leq 0$ \fg{} (proposition fausse).
  \item Pour une phrase à deux quantificateurs, on applique la règle
  successivement : la négation de \og $\forall y \in \mathbb{Z},\ \exists x \in
  \mathbb{R},\ x<y$ \fg{} est \og $\exists y \in \mathbb{Z},\ \forall x \in
  \mathbb{R},\ x \geq y$ \fg.
\end{enumerate}
\end{exemple}

\section{Propriétés fondamentales du calcul propositionnel}


Toutes les propriétés suivantes se démontrent en construisant la table de
vérité des deux membres et en vérifiant qu'ils coïncident ligne par ligne
(cette méthode est appelée \textbf{preuve par table de vérité}).

\begin{theoreme}[Lois logiques fondamentales]
Soient $P$, $Q$, $R$ des propositions (ou des prédicats, les mêmes lois
s'appliquant point par point). On a les équivalences suivantes.
\begin{enumerate}[label=(\arabic*)]
  \item \textbf{Double négation :} $\overline{\overline{P}} \iff P$.
  \item \textbf{Idempotence :} $(P\wedge P) \iff P$ \quad et \quad
  $(P\vee P) \iff P$.
  \item \textbf{Commutativité :} $(P\wedge Q)\iff(Q\wedge P)$ \quad et \quad
  $(P\vee Q)\iff(Q\vee P)$.
  \item \textbf{Associativité :} $((P\wedge Q)\wedge R)\iff(P\wedge(Q\wedge
  R))$ \quad et \quad $((P\vee Q)\vee R)\iff(P\vee(Q\vee R))$.
  \item \textbf{Distributivité :}
  $((P\wedge Q)\vee R)\iff((P\vee R)\wedge(Q\vee R))$ \quad et \quad
  $((P\vee Q)\wedge R)\iff((P\wedge R)\vee(Q\wedge R))$.
  \item \textbf{Lois de De Morgan :}
  $\overline{P\wedge Q} \iff \overline{P}\vee \overline{Q}$ \quad et \quad
  $\overline{P\vee Q} \iff \overline{P}\wedge \overline{Q}$.
  \item \textbf{Implication et disjonction :}
  $(P\Rightarrow Q) \iff \overline{P}\vee Q$.
  \item \textbf{Négation de l'implication :}
  $\overline{P\Rightarrow Q} \iff P \wedge \overline{Q}$.
  \item \textbf{Contraposée :}
  $(P\Rightarrow Q) \iff (\overline{Q}\Rightarrow \overline{P})$.
  \item \textbf{Transitivité de l'implication :}
  $\big[(P\Rightarrow Q)\wedge (Q\Rightarrow R)\big] \Rightarrow (P\Rightarrow R)$.
  \item \textbf{Caractérisation de l'équivalence :}
  $(P\Leftrightarrow Q) \iff \big[(P\Rightarrow Q)\wedge(Q\Rightarrow P)\big]$.
\end{enumerate}
\end{theoreme}

\begin{proof}[Démonstration (exemples détaillés)]
Nous détaillons deux des lois les plus utilisées ; les autres se traitent de
façon analogue (voir aussi les exercices).

\medskip
\textbf{Preuve de (6), première loi de De Morgan :}
$\overline{P\wedge Q} \iff \overline{P}\vee \overline{Q}$.

On construit la table de vérité complète des deux membres.

\begin{center}
\begin{tabular}{|c|c|c|c|c|c|c|}
\hline
$P$ & $Q$ & $P\wedge Q$ & $\overline{P\wedge Q}$ & $\overline{P}$ & $\overline{Q}$ & $\overline{P}\vee\overline{Q}$ \\
\hline
1 & 1 & 1 & 0 & 0 & 0 & 0 \\
1 & 0 & 0 & 1 & 0 & 1 & 1 \\
0 & 1 & 0 & 1 & 1 & 0 & 1 \\
0 & 0 & 0 & 1 & 1 & 1 & 1 \\
\hline
\end{tabular}
\end{center}

On constate que la colonne de $\overline{P\wedge Q}$ (4\up{e} colonne) est
identique, ligne par ligne, à celle de $\overline{P}\vee\overline{Q}$
(7\up{e} colonne). Les deux propositions ont donc toujours la même valeur de
vérité : elles sont équivalentes. $\blacksquare$

\medskip
\textbf{Preuve de (9), loi de contraposée :}
$(P\Rightarrow Q) \iff (\overline{Q}\Rightarrow \overline{P})$.

\begin{center}
\begin{tabular}{|c|c|c|c|c|c|}
\hline
$P$ & $Q$ & $P\Rightarrow Q$ & $\overline{Q}$ & $\overline{P}$ & $\overline{Q}\Rightarrow \overline{P}$ \\
\hline
1 & 1 & 1 & 0 & 0 & 1 \\
1 & 0 & 0 & 1 & 0 & 0 \\
0 & 1 & 1 & 0 & 1 & 1 \\
0 & 0 & 1 & 1 & 1 & 1 \\
\hline
\end{tabular}
\end{center}

La colonne de $P\Rightarrow Q$ coïncide exactement avec celle de
$\overline{Q}\Rightarrow\overline{P}$ : les deux propositions sont donc
équivalentes. $\blacksquare$

\medskip
\textbf{Pourquoi cette méthode est-elle valide ?} Une proposition composée
n'a que deux valeurs possibles pour chaque combinaison des valeurs de vérité de
ses composantes. Si $P$ et $Q$ sont deux propositions à $n$ variables
élémentaires, il n'y a que $2^n$ combinaisons possibles. Si les deux membres
d'une équivalence coïncident sur \emph{toutes} ces combinaisons, alors ils
prennent toujours la même valeur : ils sont équivalents par définition même de
l'équivalence logique.
\end{proof}

\begin{remarque}
La loi (7), $(P\Rightarrow Q)\iff(\overline{P}\vee Q)$, est extrêmement utile :
elle permet de \emph{transformer} toute implication en une disjonction, ce qui
explique aussi pourquoi $P\Rightarrow Q$ est automatiquement vraie dès que $P$
est fausse (car alors $\overline{P}$ est vraie, donc $\overline{P}\vee Q$ l'est
aussi, quelle que soit la valeur de $Q$).
\end{remarque}


\section{Les grands types de raisonnement}


\subsection{Le raisonnement déductif}

\begin{proposition}[Principe]
Pour montrer qu'une proposition $Q$ est vraie, il suffit de montrer que les
deux propositions $P\Rightarrow Q$ et $P$ sont vraies (pour une proposition
$P$ bien choisie).
\end{proposition}

Ce raisonnement repose sur le fait que $\big[P \wedge (P\Rightarrow
Q)\big] \Rightarrow Q$ est \textbf{toujours vraie}, comme le montre sa table
de vérité :

\begin{center}
\begin{tabular}{|c|c|c|c|c|}
\hline
$P$ & $Q$ & $P\Rightarrow Q$ & $P\wedge(P\Rightarrow Q)$ & $[P\wedge(P\Rightarrow Q)]\Rightarrow Q$\\
\hline
1 & 1 & 1 & 1 & 1 \\
1 & 0 & 0 & 0 & 1 \\
0 & 1 & 1 & 0 & 1 \\
0 & 0 & 1 & 0 & 1 \\
\hline
\end{tabular}
\end{center}

La dernière colonne ne contient que des $1$ : la proposition est une
\textbf{tautologie} (toujours vraie, quelles que soient les valeurs de $P$ et
$Q$).

\begin{exemple}
Montrons que l'équation $x^2-x+3=0$ n'a pas de solution réelle.

On sait ($P$) que pour une équation du second degré $ax^2+bx+c=0$ ($a\neq0$)
dans $\mathbb{R}$, on a l'implication vraie : \og $\Delta<0$ $\Rightarrow$
l'équation n'a pas de solution réelle \fg{} ($Q$ dépend ici de $\Delta$).

Ici $a=1$, $b=-1$, $c=3$, donc $\Delta = b^2-4ac = 1-12=-11<0$. La proposition
$P$ (\og $\Delta<0$ \fg) est donc vraie. Par raisonnement déductif, on conclut
que $Q$ est vraie : l'équation n'a pas de solution réelle.
\end{exemple}

\subsection{Le raisonnement par contraposée}

\begin{proposition}[Principe]
Pour montrer que $P\Rightarrow Q$ est vraie, il suffit de montrer que sa
contraposée $\overline{Q}\Rightarrow \overline{P}$ est vraie.
\end{proposition}

Ce raisonnement est justifié par la loi (9) du Théorème~1.1 : $P\Rightarrow Q$
et $\overline{Q}\Rightarrow\overline{P}$ sont toujours équivalentes. On
utilise ce mode de raisonnement quand il est plus facile de manipuler
$\overline{Q}$ (l'hypothèse) que $P$ directement.

\begin{exemple}
Montrons que, pour tout entier $a$ : (\og $a^2$ est pair \fg) $\Rightarrow$
(\og $a$ est pair \fg).

Il est plus simple de montrer la contraposée : (\og $a$ n'est pas pair \fg)
$\Rightarrow$ (\og $a^2$ n'est pas pair \fg), c'est-à-dire (\og $a$ est
impair \fg) $\Rightarrow$ (\og $a^2$ est impair \fg).

Supposons $a$ impair. Alors il existe un entier $n$ tel que $a=2n+1$. Calculons
$a^2$ :
\[
a^2 = (2n+1)^2 = 4n^2+4n+1 = 2(2n^2+2n)+1.
\]
En posant $k=2n^2+2n\in\mathbb{Z}$, on obtient $a^2=2k+1$, qui est bien un
nombre impair. La contraposée est donc démontrée, ce qui prouve, par la loi de
contraposée, l'implication initiale.
\end{exemple}

\subsection{Le raisonnement par l'absurde}

\begin{proposition}[Principe]
Pour montrer qu'une proposition $Q$ est vraie, on suppose que sa négation
$\overline{Q}$ est vraie, et on montre que cette supposition conduit à une
contradiction (une proposition fausse, ou deux propositions contradictoires).
On en déduit que $\overline{Q}$ est fausse, donc que $Q$ est vraie.
\end{proposition}

\begin{exemple}
Montrons que $\sqrt{2}$ est irrationnel (c'est-à-dire $\sqrt2 \notin
\mathbb{Q}$).

Supposons, par l'absurde, que $\sqrt{2}\in\mathbb{Q}$. Alors on peut écrire
$\sqrt2 = \dfrac{a}{b}$, où $a$ et $b$ sont deux entiers premiers entre eux
(la fraction est sous forme irréductible) et $b\neq 0$.

En élevant au carré : $2=\dfrac{a^2}{b^2}$, donc $a^2=2b^2$. Ainsi $a^2$ est
pair. D'après l'exemple~1.13, cela entraîne que $a$ est pair : il existe donc
un entier $a'$ tel que $a=2a'$.

En remplaçant : $(2a')^2=2b^2$, soit $4a'^2=2b^2$, c'est-à-dire
$b^2=2a'^2$. Donc $b^2$ est pair, et par le même argument, $b$ est pair.

Ainsi $a$ et $b$ sont tous deux pairs : ils sont divisibles par $2$. Mais nous
avions supposé $a$ et $b$ premiers entre eux (aucun diviseur commun autre que
$1$) : \textbf{contradiction}.

Cette contradiction montre que l'hypothèse de départ (\og $\sqrt2$ est
rationnel \fg) est fausse. Donc $\sqrt2$ est irrationnel.
\end{exemple}

\subsection{Le raisonnement par disjonction de cas}

\begin{proposition}[Principe]
Pour montrer qu'une proposition $Q$ est vraie, il suffit de trouver une
proposition $P$ telle que les deux implications $P\Rightarrow Q$ et
$\overline{P}\Rightarrow Q$ soient vraies. On étudie ainsi \emph{tous les cas
possibles} (le cas où $P$ est vraie, et le cas où $P$ est fausse), qui
couvrent toutes les situations.
\end{proposition}

\begin{exemple}
Montrons que, pour tout entier $n$, le produit $n(n+1)$ est pair.

\textbf{1\up{er} cas :} $n$ est pair. Alors il existe $k\in\mathbb{Z}$ tel que
$n=2k$. On a $n(n+1) = 2k(2k+1)$, qui est bien pair (multiple de $2$).

\textbf{2\up{e} cas :} $n$ est impair. Alors il existe $k\in\mathbb{Z}$ tel que
$n=2k+1$. On a $n(n+1) = (2k+1)(2k+2) = 2(2k+1)(k+1)$, qui est bien pair.

Dans les deux cas (qui couvrent toutes les possibilités pour un entier $n$),
$n(n+1)$ est pair. La proposition est donc démontrée pour tout entier $n$.
\end{exemple}

\subsection{Le raisonnement par récurrence}


\begin{proposition}[Principe de récurrence]
Soit $P(n)$ un prédicat défini pour tout entier naturel $n \geq n_0$. Si
\begin{enumerate}
  \item (\textbf{Initialisation}) $P(n_0)$ est vraie, et
  \item (\textbf{Hérédité}) pour tout $n\geq n_0$, $P(n)\Rightarrow P(n+1)$
  est vraie,
\end{enumerate}
alors $P(n)$ est vraie pour tout entier $n\geq n_0$.
\end{proposition}

\begin{remarque}[Intuition : l'effet domino]
On peut se représenter ce principe comme une rangée infinie de dominos
numérotés $n_0, n_0+1, n_0+2,\dots$ L'hérédité garantit que si un domino tombe,
il fait tomber le suivant. L'initialisation garantit que le premier domino
tombe effectivement. La conclusion : \emph{tous} les dominos tombent, l'un
après l'autre, sans exception.
\end{remarque}

\begin{exemple}
Montrons par récurrence que pour tout entier $n\geq 1$ :
\[
1+2+\cdots+n = \frac{n(n+1)}{2}.
\]
Notons $P(n)$ cette égalité.

\textbf{Initialisation ($n=1$) :} Le membre de gauche vaut $1$. Le membre de
droite vaut $\dfrac{1\times 2}{2}=1$. Les deux membres sont égaux : $P(1)$ est
vraie.

\textbf{Hérédité :} Supposons que $P(n)$ soit vraie pour un entier $n\geq 1$
fixé (c'est l'\textbf{hypothèse de récurrence}), c'est-à-dire
$1+2+\cdots+n=\dfrac{n(n+1)}{2}$. Montrons que $P(n+1)$ est vraie, c'est-à-dire
que $1+2+\cdots+n+(n+1) = \dfrac{(n+1)(n+2)}{2}$.

Partons du membre de gauche de $P(n+1)$ et utilisons l'hypothèse de
récurrence :
\[
1+2+\cdots+n+(n+1) = \underbrace{\big(1+2+\cdots+n\big)}_{=\,\frac{n(n+1)}{2}\ \text{par hypothèse}} + (n+1)
= \frac{n(n+1)}{2}+(n+1).
\]
On met $(n+1)$ en facteur :
\[
\frac{n(n+1)}{2}+(n+1) = (n+1)\left(\frac{n}{2}+1\right) = (n+1)\cdot\frac{n+2}{2} = \frac{(n+1)(n+2)}{2}.
\]
On obtient exactement le membre de droite attendu pour $P(n+1)$. L'hérédité
est donc démontrée.

\textbf{Conclusion :} $P(1)$ est vraie et pour tout $n\geq1$, $P(n)\Rightarrow
P(n+1)$. D'après le principe de récurrence, $P(n)$ est vraie pour tout entier
$n\geq 1$.
\end{exemple}


\section{Illustration synthétique}

Le schéma logique ci-dessous résume l'articulation des notions du chapitre :
des propositions de base, on construit des propositions composées à l'aide des
connecteurs ; en introduisant des variables, on passe aux prédicats, que l'on
transforme en propositions à l'aide des quantificateurs ; enfin, ces outils
permettent de bâtir des démonstrations selon les cinq modes de raisonnement.

\begin{center}
\begin{tabular}{|>{\centering\arraybackslash}p{4.2cm}|>{\centering\arraybackslash}p{4.2cm}|>{\centering\arraybackslash}p{4.2cm}|}
\hline
\textbf{Objets de base} & \textbf{Outils de construction} & \textbf{Objets construits}\\
\hline
Propositions $P,Q,R$ & $\overline{\ \cdot\ },\ \wedge,\ \vee,\ \Rightarrow,\ \Leftrightarrow$ & Propositions composées\\
\hline
Prédicats $P(x)$ & $\forall,\ \exists$ & Propositions quantifiées\\
\hline
Propositions vraies établies & Raisonnement (déductif, contraposée, absurde, cas, récurrence) & Nouveaux théorèmes\\
\hline
\end{tabular}
\end{center}


\chapter{Théorie ensembliste et relations}

\section{Introduction}


Ce chapitre a pour objectifs :
\begin{itemize}
  \item de définir rigoureusement un \textbf{ensemble}, ses
  \textbf{éléments}, et la relation d'\textbf{appartenance} ;
  \item de manipuler les opérations sur les ensembles : \textbf{inclusion},
  \textbf{intersection}, \textbf{union}, \textbf{différence},
  \textbf{complémentaire} ;
  \item d'introduire le \textbf{produit cartésien} de deux ensembles ;
  \item de définir une \textbf{relation binaire} et ses propriétés
  fondamentales (réflexivité, symétrie, antisymétrie, transitivité) ;
  \item d'étudier en détail les \textbf{relations d'équivalence} (classes,
  ensembles quotients) et les \textbf{relations d'ordre} (ordre total,
  ordre partiel).
\end{itemize}

Ce chapitre s'appuie directement sur le Chapitre 1 : les prédicats
$P(x)$ serviront à décrire les ensembles \og en compréhension \fg, et les
connecteurs logiques ($\wedge$, $\vee$, $\Rightarrow$) serviront à définir les
opérations ensemblistes et les propriétés des relations. Il prépare le
Chapitre 3, consacré aux applications (fonctions), qui sont un cas particulier
de relations.

\begin{application}
La théorie des ensembles est le langage universel des mathématiques : elle
sert de fondation à l'algèbre, à l'analyse, à la géométrie et aux
probabilités. En \textbf{informatique}, les ensembles modélisent les types
\texttt{set} des langages de programmation ; les relations d'équivalence
sont à la base du \textbf{partitionnement de données} (clustering) et les
relations d'ordre gouvernent tous les \textbf{algorithmes de tri}.
\end{application}

\section{Prérequis}

Ce chapitre utilise le vocabulaire logique du Chapitre 1 : propositions,
prédicats, connecteurs $\overline{\ \cdot\ }, \wedge, \vee, \Rightarrow,
\Leftrightarrow$, quantificateurs $\forall$, $\exists$. Il est indispensable
de maîtriser en particulier la négation d'une proposition quantifiée avant
d'aborder la section sur les relations d'ordre.

\section{Ensembles}

\subsection{Notion d'ensemble}

\begin{definition}
Un \textbf{ensemble} est une collection d'objets, appelés \textbf{éléments}
de cet ensemble. Si un ensemble est noté $A$ :
\begin{itemize}
  \item on écrit $x \in A$ pour dire que $x$ est un élément de $A$ (\og $x$
  appartient à $A$ \fg) ;
  \item on écrit $x \notin A$ pour dire que $x$ n'est pas un élément de $A$.
\end{itemize}
\end{definition}

Un ensemble peut être décrit de deux façons :

\begin{definition}[Description en extension]
On décrit un ensemble $A$ \textbf{en extension} en donnant la liste complète
de ses éléments : $A = \{a_1, a_2, \dots, a_n\}$.
\end{definition}

\begin{definition}[Description en compréhension]
On décrit un ensemble $A$ \textbf{en compréhension} en caractérisant ses
éléments parmi ceux d'un ensemble $E$ déjà connu, au moyen d'un prédicat
$P(x)$ : $A = \{x \in E : P(x)\}$ (on lit : \og l'ensemble des $x$ de $E$
tels que $P(x)$ \fg).
\end{definition}

\begin{exemple}
\begin{enumerate}
  \item $F = \{0,1,2,3,4,5,6,7,8,9\}$ : ensemble des chiffres du système
  décimal (description en extension).
  \item $2\mathbb{N} = \{x \in \mathbb{N} : 2 \text{ divise } x\}$ : ensemble
  des entiers naturels pairs (description en compréhension).
  \item $\mathbb{R}^+ = \{x \in \mathbb{R} : x \geq 0\}$ : ensemble des réels
  positifs.
  \item On a $3 \in F$, $7 \notin 2\mathbb{N}$, $18 \in 2\mathbb{N}$,
  $-4 \notin \mathbb{R}^+$.
\end{enumerate}
\end{exemple}

\begin{definition}[2.4 -- Ensemble vide]
L'ensemble qui ne contient aucun élément est appelé \textbf{ensemble vide},
noté $\varnothing$.
\end{definition}

\subsection{Inclusion et parties d'un ensemble}


\begin{definition}[2.5 -- Inclusion]
Soient $A$ et $B$ deux ensembles. On dit que $B$ est un \textbf{sous-ensemble}
(ou une \textbf{partie}) de $A$ si tout élément de $B$ est un élément de $A$ :
\[
B \subset A \quad \Longleftrightarrow \quad \forall x,\ (x \in B \Rightarrow x \in A).
\]
Dans le cas contraire, on écrit $B \not\subset A$.
\end{definition}

\begin{remarque}
Par convention, $\varnothing \subset E$ pour tout ensemble $E$ : cette
affirmation est \emph{vraie} car son hypothèse (\og $x \in \varnothing$ \fg)
est toujours fausse, et une implication à hypothèse fausse est vraie (voir
Chapitre 1, section 1.3.5).
\end{remarque}


\begin{definition}[Ensemble des parties]
L'ensemble de toutes les parties d'un ensemble $A$ est noté $\mathcal{P}(A)$.
\end{definition}

\begin{exemple}
Soit $A=\{0,1,2\}$. Alors
\[
\mathcal{P}(A) = \big\{\varnothing,\ \{0\},\ \{1\},\ \{2\},\ \{0,1\},\ \{0,2\},\ \{1,2\},\ A\big\}.
\]
On remarque que $\mathcal{P}(A)$ possède $8=2^3$ éléments : c'est un fait
général, si $A$ possède $n$ éléments, $\mathcal{P}(A)$ en possède $2^n$
(chaque élément de $A$ pouvant être \og inclus \fg{} ou \og exclu \fg{} d'une
partie, indépendamment des autres).
\end{exemple}

\subsection{Opérations sur les ensembles}


\begin{definition}
Soient $E$, $A$ et $B$ des ensembles.
\begin{enumerate}
  \item La \textbf{différence} de $A$ et $B$ est
  $A\setminus B = \{x \in A : x \notin B\}$.
  \item Si $A \subset E$, le \textbf{complémentaire} de $A$ dans $E$ est
  $E\setminus A$, noté $C_E A$ (ou $A^c$) : $C_E A = \{x \in E : x \notin A\}$.
  \item L'\textbf{intersection} de $A$ et $B$ est
  $A \cap B = \{x : x \in A \text{ et } x \in B\}$.
  \item L'\textbf{union} (ou réunion) de $A$ et $B$ est
  $A \cup B = \{x : x \in A \text{ ou } x \in B\}$.
\end{enumerate}
\end{definition}

\begin{exemple}
Soit $F=\{0,1,2,3,4,5,6,7,8,9\}$.
\begin{enumerate}
  \item $F \subset \mathbb{N}$ mais $F \not\subset 2\mathbb{N}$ (car, par
  exemple, $1 \in F$ mais $1 \notin 2\mathbb{N}$).
  \item $F \setminus 2\mathbb{N} = \{1,3,5,7,9\}$ (les éléments de $F$ qui
  \emph{ne sont pas} pairs).
  \item $C_{\mathbb{N}} F = \{x \in \mathbb{N} : x \geq 10\} =
  \{10,11,12,\dots\}$.
  \item $F \cap 2\mathbb{N} = \{0,2,4,6,8\}$ (les chiffres pairs).
  \item $\{0,2,4,6,8\} \cup \{1,2,4\} = \{0,1,2,4,6,8\}$.
\end{enumerate}
\end{exemple}

\subsubsection*{Illustration : diagrammes de Venn}

Les opérations ensemblistes se visualisent naturellement à l'aide de
\textbf{diagrammes de Venn}, où chaque ensemble est représenté par une région
du plan (souvent un disque), et les opérations correspondent à des régions
du dessin.

\begin{center}
\begin{figure}[h]
  \centering
  \includegraphics[width=0.6\textwidth]{image1}
  %\caption{Graphe de la fonction $f$}
\end{figure}
\hspace{1cm}
\begin{figure}[h]
  \centering
  \includegraphics[width=0.6\textwidth]{image2}
  %\caption{Graphe de la fonction $f$}
\end{figure}
\end{center}

\begin{proposition}[Propriétés des opérations ensemblistes]
Soient $A$, $B$, $C$ des parties d'un ensemble $E$. On a :
\begin{enumerate}[label=(\arabic*)]
  \item $A \cap A = A$ et $A \cup A = A$ (idempotence) ;
  \item $A \cap B = B \cap A$ et $A \cup B = B \cup A$ (commutativité) ;
  \item $(A\cap B)\cap C = A \cap (B\cap C)$ et $(A\cup B)\cup C = A \cup
  (B\cup C)$ (associativité) ;
  \item $(A\cap B)\cup C = (A\cup C)\cap(B\cup C)$ et
  $(A\cup B)\cap C = (A\cap C)\cup(B\cap C)$ (distributivité) ;
  \item Si $A,B \subset E$ : $C_E(A \cap B) = (C_E A) \cup (C_E B)$ et
  $C_E(A \cup B) = (C_E A) \cap (C_E B)$ (lois de De Morgan ensemblistes).
\end{enumerate}
\end{proposition}

\begin{proof}
Ces propriétés se démontrent en traduisant les opérations ensemblistes en
propositions logiques sur les prédicats correspondants, puis en appliquant
les lois logiques du Chapitre 1 (Théorème 1.1). Détaillons la loi de De
Morgan $C_E(A\cap B) = (C_E A)\cup(C_E B)$.

Soit $A=\{x\in E : P(x)\}$ et $B=\{x\in E : Q(x)\}$, où $P$ et $Q$ sont les
prédicats caractérisant $A$ et $B$ dans $E$. Alors, pour tout $x \in E$ :
\[
x \in C_E(A\cap B)
\iff x \notin (A \cap B)
\iff \overline{P(x) \wedge Q(x)}
\iff \overline{P(x)} \vee \overline{Q(x)} \quad \text{(De Morgan logique, Th. 1.1)}
\]
\[
\iff (x \notin A) \vee (x \notin B)
\iff x \in C_E A \ \text{ ou } \ x \in C_E B
\iff x \in (C_E A) \cup (C_E B).
\]
Les deux ensembles $C_E(A\cap B)$ et $(C_E A)\cup(C_E B)$ contiennent donc
exactement les mêmes éléments : ils sont égaux. La seconde loi de De Morgan
se démontre de façon analogue en utilisant l'autre loi de De Morgan
logique.
\end{proof}

\begin{remarque}
Cette démonstration illustre un principe fondamental : \textbf{toute
propriété ensembliste se ramène, via les descriptions en compréhension, à une
propriété logique déjà connue}. C'est pourquoi le Chapitre 1 était un
prérequis indispensable.
\end{remarque}

\subsection{Partition d'un ensemble}


\begin{definition}
On appelle \textbf{partition} d'un ensemble $E$ toute famille $\mathcal{F}$ de
parties \textbf{non vides} de $E$ telle que :
\begin{enumerate}
  \item les éléments de $\mathcal{F}$ sont deux à deux \textbf{disjoints} :
  $\forall A,B \in \mathcal{F},\ A\neq B \Rightarrow A\cap B = \varnothing$ ;
  \item $\mathcal{F}$ \textbf{recouvre} $E$ : $\displaystyle\bigcup_{A \in
  \mathcal{F}} A = E$.
\end{enumerate}
\end{definition}

\begin{exemple}
Soit $B=\{0,1,2,3,4\}$. La famille $\mathcal{F} = \big\{\{0,2\}, \{1,4\},
\{3\}\big\}$ est une partition de $B$ : les trois parties sont non vides,
deux à deux disjointes, et leur réunion $\{0,2\}\cup\{1,4\}\cup\{3\}=B$.
En revanche, $\mathcal{F}' = \big\{\{0,1\}, \{3,4\}, \{1,2\}\big\}$
\textbf{n'est pas} une partition de $B$ : $\{0,1\}$ et $\{1,2\}$ ne sont pas
disjoints (elles partagent l'élément $1$), et $0$ n'apparaît nulle part ne
pose pas problème ici mais la non-disjonction suffit à invalider la
partition.
\end{exemple}

\section{Produit cartésien}


\begin{definition}
Le \textbf{produit cartésien} de deux ensembles $A$ et $B$, noté $A \times
B$, est l'ensemble des \textbf{couples} $(a,b)$ tels que $a \in A$ et
$b \in B$ :
\[
A \times B = \{(a,b) : a \in A \text{ et } b \in B\}.
\]
\end{definition}

\begin{remarque}
Un couple $(a,b)$ est un objet \textbf{ordonné} : $(a,b) \neq (b,a)$ dès que
$a \neq b$. C'est ce qui le distingue d'un ensemble $\{a,b\}$, où l'ordre
d'écriture n'a pas d'importance ($\{a,b\}=\{b,a\}$).
\end{remarque}

\begin{exemple}
Soit $A=\{0,2,4\}$ et $B=\{1,2\}$. Alors
\[
A \times B = \big\{(0,1),(0,2),(2,1),(2,2),(4,1),(4,2)\big\},
\]
tandis que
\[
B \times A = \big\{(1,0),(2,0),(1,2),(2,2),(1,4),(2,4)\big\}.
\]
On constate que $A \times B \neq B \times A$ en général (par exemple
$(0,1) \in A\times B$ mais $(0,1) \notin B\times A$) : le produit cartésien
\textbf{n'est pas commutatif}.
\end{exemple}


\begin{application}
Le produit cartésien $\mathbb{R} \times \mathbb{R} = \mathbb{R}^2$ est
précisément l'ensemble des coordonnées $(x,y)$ du plan géométrique : c'est le
pont entre l'algèbre et la géométrie (géométrie analytique de Descartes, d'où
le nom \og cartésien \fg). En informatique, le produit cartésien de deux
tables est une opération fondamentale des bases de données relationnelles
(\texttt{CROSS JOIN} en SQL).
\end{application}

\section{Relations binaires sur un ensemble}

\subsection{Notion de relation binaire}


\begin{definition}
On appelle \textbf{relation} d'un ensemble $A$ vers un ensemble $B$ toute
correspondance $\mathcal{R}$ qui lie certains éléments de $A$ à certains
éléments de $B$. Si $x$ est lié à $y$ par $\mathcal{R}$, on écrit
$x\mathcal{R}y$ (sinon $x\cancel{\mathcal{R}}y$). Une relation de $A$ vers
$A$ est dite \textbf{relation sur $A$}.
\end{definition}

\begin{definition}[Graphe d'une relation]
Le \textbf{graphe} de la relation $\mathcal{R}$ (de $A$ vers $B$) est
l'ensemble
\[
G_{\mathcal{R}} = \{(x,y) \in A\times B : x\,\mathcal{R}\,y\}.
\]
Une relation est entièrement déterminée par son graphe : on identifie souvent
$\mathcal{R}$ à une partie de $A \times B$.
\end{definition}

\begin{exemple}[2.6]
La relation de \textbf{divisibilité} $\mathcal{R}_d$ sur $\mathbb{Z}$, définie
par $x\,\mathcal{R}_d\,y \iff x$ divise $y$, vérifie par exemple
$1\,\mathcal{R}_d\,x$ et $x\,\mathcal{R}_d\,0$ pour tout $x \in \mathbb{Z}$
(car $1$ divise tout entier, et tout entier divise $0$). On a
$3 \,\mathcal{R}_d\, (-21)$ (car $-21=3\times(-7)$), mais
$3\,\cancel{\mathcal{R}_d}\,20$.
\end{exemple}

\subsection{Propriétés d'une relation binaire}


\begin{definition}
Soit $\mathcal{R}$ une relation sur un ensemble $A$.
\begin{enumerate}
  \item $\mathcal{R}$ est \textbf{réflexive} si $\forall x \in A,\ x\mathcal{R}x$.
  \item $\mathcal{R}$ est \textbf{symétrique} si $\forall x,y \in A,\
  x\mathcal{R}y \Rightarrow y\mathcal{R}x$.
  \item $\mathcal{R}$ est \textbf{antisymétrique} si $\forall x,y \in A,\
  (x\mathcal{R}y \wedge y\mathcal{R}x) \Rightarrow x=y$.
  \item $\mathcal{R}$ est \textbf{transitive} si $\forall x,y,z \in A,\
  (x\mathcal{R}y \wedge y\mathcal{R}z) \Rightarrow x\mathcal{R}z$.
\end{enumerate}
\end{definition}


\begin{exemple}
Étudions la relation de divisibilité $\mathcal{R}_d$ sur $\mathbb{Z}$ :
$x\,\mathcal{R}_d\,y \iff x \text{ divise } y$.

\textbf{Réflexivité :} pour tout $x\in\mathbb{Z}$, $x$ divise $x$ (car
$x = 1\times x$). Donc $\mathcal{R}_d$ est réflexive.

\textbf{Symétrie :} $1$ divise $4$, mais $4$ ne divise pas $1$. Donc
$\mathcal{R}_d$ n'est \textbf{pas} symétrique.

\textbf{Antisymétrie :} $1$ divise $-1$ et $-1$ divise $1$, mais $1 \neq -1$.
Donc $\mathcal{R}_d$ n'est \textbf{pas} antisymétrique (sur $\mathbb{Z}$ tout
entier ; nous verrons en section 2.6 qu'elle le devient si l'on se restreint
à $\mathbb{N}$).

\textbf{Transitivité :} si $x$ divise $y$ et $y$ divise $z$, alors il existe
$k,k'\in\mathbb{Z}$ tels que $y=kx$ et $z=k'y=k'kx$, donc $x$ divise $z$.
$\mathcal{R}_d$ est transitive.
\end{exemple}

\section{Relation d'équivalence, classes d'équivalence et ensemble quotient}


\begin{definition}
Une relation $\mathcal{R}$ sur un ensemble $A$ est une \textbf{relation
d'équivalence} si elle est \textbf{réflexive}, \textbf{symétrique} et
\textbf{transitive}.
\end{definition}

\begin{definition}[Classe d'équivalence, ensemble quotient]
Soit $\mathcal{R}$ une relation d'équivalence sur $A$.
\begin{enumerate}
  \item Pour $a \in A$, la \textbf{classe d'équivalence} de $a$ modulo
  $\mathcal{R}$ est $\dot{a} = \{x \in A : x\,\mathcal{R}\,a\}$.
  \item L'ensemble $A/\mathcal{R} = \{\dot{a} : a \in A\}$ est appelé
  \textbf{ensemble quotient} de $A$ par $\mathcal{R}$.
\end{enumerate}
\end{definition}

\begin{theoreme}[Les classes d'équivalence forment une partition]
Si $\mathcal{R}$ est une relation d'équivalence sur $A$, l'ensemble des
classes d'équivalence distinctes forme une \textbf{partition} de $A$.
\end{theoreme}

\begin{proof}
Il faut vérifier les deux conditions de la Définition 2.8.

\textbf{(i) Chaque classe est non vide et leur réunion vaut $A$.} Soit $a\in
A$. Par réflexivité, $a\,\mathcal{R}\,a$, donc $a \in \dot a$ : chaque classe
$\dot a$ est non vide (elle contient au moins $a$), et tout élément $a$ de
$A$ appartient à sa propre classe, donc $\bigcup_{a\in A} \dot a = A$.

\textbf{(ii) Deux classes distinctes sont disjointes.} Montrons la
contraposée : si $\dot a \cap \dot b \neq \varnothing$, alors $\dot a =
\dot b$. Soit $c \in \dot a \cap \dot b$, c'est-à-dire $c\,\mathcal{R}\,a$ et
$c\,\mathcal{R}\,b$. Par symétrie, $a\,\mathcal{R}\,c$ ; par transitivité avec
$c\,\mathcal{R}\,b$, on obtient $a\,\mathcal{R}\,b$.

Montrons alors $\dot a \subset \dot b$ : soit $x \in \dot a$, c'est-à-dire
$x\,\mathcal{R}\,a$. Comme $a\,\mathcal{R}\,b$, la transitivité donne
$x\,\mathcal{R}\,b$, donc $x \in \dot b$. L'inclusion réciproque $\dot b
\subset \dot a$ s'obtient de façon symétrique (en échangeant les rôles de $a$
et $b$, licite car $a\,\mathcal{R}\,b \Rightarrow b\,\mathcal{R}\,a$ par
symétrie). Donc $\dot a = \dot b$.

Ainsi, deux classes qui se rencontrent sont en fait égales : des classes
\emph{distinctes} sont donc nécessairement \emph{disjointes}. Les deux
conditions d'une partition sont vérifiées.
\end{proof}

\begin{remarque}[Pourquoi chaque hypothèse est nécessaire]
La réflexivité garantit que chaque élément appartient bien à \emph{une}
classe (recouvrement). La symétrie et la transitivité sont toutes deux
utilisées pour montrer que deux classes qui se rencontrent sont confondues :
sans l'une ou l'autre, on pourrait avoir des classes qui se chevauchent
partiellement sans être égales, ce qui romprait la structure de partition.
\end{remarque}

\begin{exemple}
Soit $\mathcal{R}_e$ la relation de \textbf{congruence modulo $n$}
(où $n \in \mathbb{N}^*$ est fixé) sur $\mathbb{Z}$, définie par
\[
x\,\mathcal{R}_e\,y \iff n \text{ divise } x-y.
\]

\textbf{Vérifions que $\mathcal{R}_e$ est une relation d'équivalence.}

\textit{Réflexivité :} pour tout $x$, $x-x=0=n\times 0$, donc $n$ divise
$x-x$ : $x\,\mathcal{R}_e\,x$.

\textit{Symétrie :} si $n$ divise $x-y$, alors $x-y=nk$ pour un
$k\in\mathbb{Z}$, donc $y-x = n(-k)$, et $n$ divise aussi $y-x$ :
$y\,\mathcal{R}_e\,x$.

\textit{Transitivité :} si $n$ divise $x-y$ et $n$ divise $y-z$, alors
$x-y=nk$ et $y-z=nk'$ pour $k,k'\in\mathbb{Z}$, donc
$x-z=(x-y)+(y-z)=n(k+k')$ : $n$ divise $x-z$, donc $x\,\mathcal{R}_e\,z$.

$\mathcal{R}_e$ est donc bien une relation d'équivalence.

\textbf{Classes d'équivalence.} Pour $a \in \mathbb{Z}$ :
\[
\dot a = \{x \in \mathbb{Z} : n \text{ divise } x-a\} = \{x \in \mathbb{Z} :
x = a+nq,\ q\in\mathbb{Z}\} = n\mathbb{Z}+a.
\]
On montre (division euclidienne de $a$ par $n$) que ces classes ne dépendent,
en réalité, que du \textbf{reste} de $a$ dans la division par $n$, et qu'il y
en a exactement $n$ distinctes :
\[
\mathbb{Z}/\mathcal{R}_e = \{\dot 0, \dot 1, \dots, \dot{n-1}\}.
\]
Cet ensemble quotient est noté $\mathbb{Z}/n\mathbb{Z}$ : il joue un rôle
central dans toute la suite du cours d'algèbre (structures de groupe,
d'anneau, de corps).
\end{exemple}

\begin{application}
La congruence modulo $n$ est utilisée en \textbf{cryptographie}
(algorithme RSA), dans le calcul des \textbf{clés de contrôle} (numéro de
sécurité sociale, IBAN), et dans la conception des \textbf{tables de
hachage} en informatique, où l'on regroupe des données selon le reste modulo
la taille de la table.
\end{application}

\section{Relation d'ordre et ensembles ordonnés}


\begin{definition}
Une relation $\mathcal{R}$ sur un ensemble $A$ est une \textbf{relation
d'ordre} si elle est \textbf{réflexive}, \textbf{antisymétrique} et
\textbf{transitive}. On note alors souvent $\leq_{\mathcal{R}}$ au lieu de
$\mathcal{R}$.
\end{definition}

\begin{definition}[ Ordre total, ordre partiel]
Soit $\leq_{\mathcal{R}}$ une relation d'ordre sur $A$.
\begin{itemize}
  \item L'ordre est \textbf{total} si $\forall x,y \in A,\ (x
  \leq_{\mathcal{R}} y) \vee (y \leq_{\mathcal{R}} x)$ (deux éléments
  quelconques sont toujours comparables).
  \item L'ordre est \textbf{partiel} si $\exists x,y \in A,\
  (x\not\leq_{\mathcal{R}} y) \wedge (y \not\leq_{\mathcal{R}} x)$ (il existe
  au moins deux éléments incomparables).
\end{itemize}
Deux éléments $x$ et $y$ sont dits \textbf{comparables} par
$\leq_{\mathcal{R}}$ si $x\leq_{\mathcal{R}} y$ ou $y\leq_{\mathcal{R}} x$.
\end{definition}

\begin{exemple}
\begin{enumerate}
  \item La relation usuelle $\leq$ sur $\mathbb{R}$ est un \textbf{ordre
  total} : deux réels quelconques sont toujours comparables.
  \item La relation de divisibilité $\mathcal{R}_d$ n'est \textbf{pas} une
  relation d'ordre sur $\mathbb{Z}$ (elle n'est pas antisymétrique, voir
  Exemple 2.7), mais elle en devient une si l'on se restreint à
  $\mathbb{N}$ : c'est un \textbf{ordre partiel} sur $\mathbb{N}$, car par
  exemple $3$ et $7$ ne sont pas comparables ($3$ ne divise pas $7$ et $7$ ne
  divise pas $3$).
  \item L'inclusion $\subset$ sur $\mathcal{P}(E)$ (l'ensemble des parties
  d'un ensemble $E$ ayant au moins deux éléments) est un \textbf{ordre
  partiel} : si $E=\{1,2\}$, les parties $\{1\}$ et $\{2\}$ ne sont pas
  comparables par $\subset$.
\end{enumerate}
\end{exemple}

\begin{theoreme}
La relation de divisibilité, restreinte à $\mathbb{N}$, est une relation
d'ordre (partiel) sur $\mathbb{N}$.
\end{theoreme}

\begin{proof}
La réflexivité et la transitivité ont déjà été établies dans l'Exemple 2.7
(elles restent valables sur $\mathbb{N}$). Il reste à montrer
l'\textbf{antisymétrie} sur $\mathbb{N}$ (elle était mise en défaut sur
$\mathbb{Z}$ à cause des signes).

Soient $a,b \in \mathbb{N}$ tels que $a$ divise $b$ et $b$ divise $a$. Il
existe donc $q,q' \in \mathbb{N}$ tels que $b=qa$ et $a=q'b$. En substituant :
\[
a = q'b = q'(qa) = (qq')a.
\]
\textbf{Cas $a=0$ :} alors $b=qa=0$, donc $a=b=0$ : l'égalité est vérifiée.

\textbf{Cas $a\neq 0$ :} on peut simplifier par $a$ dans $a=(qq')a$, ce qui
donne $qq'=1$. Comme $q,q'\in\mathbb{N}$, la seule possibilité est $q=q'=1$
(un produit de deux entiers naturels ne vaut $1$ que si les deux facteurs
valent $1$). Donc $b=qa=a$.

Dans les deux cas, $a=b$ : la relation de divisibilité est bien antisymétrique
sur $\mathbb{N}$, donc c'est une relation d'ordre. Cet ordre n'est pas total
(exemple 2.9), c'est donc un ordre partiel.
\end{proof}

\begin{remarque}[Pourquoi la restriction à $\mathbb{N}$ est nécessaire]
La démonstration précédente échoue sur $\mathbb{Z}$ : l'équation $qq'=1$ avec
$q,q'\in\mathbb{Z}$ admet \emph{deux} solutions, $(q,q')=(1,1)$ mais aussi
$(q,q')=(-1,-1)$, qui correspond au cas $a=-b$ avec $a\neq b$ dès que $a\neq
0$ (voir le contre-exemple $1$ divise $-1$ et $-1$ divise $1$, mais $1\neq
-1$, de l'exemple 2.7). C'est exactement pourquoi une hypothèse
(\og $\mathbb{N}$ \fg{} plutôt que \og $\mathbb{Z}$ \fg) est indispensable
pour ce théorème.
\end{remarque}

\begin{exemple}[2.10]
La façon dont les mots sont rangés dans un dictionnaire définit une relation
d'\textbf{ordre total} sur l'ensemble des mots, appelée \textbf{ordre
lexicographique}, noté $\leq_{\text{lex}}$. On a par exemple
$\text{alg\`ebre} \leq_{\text{lex}} \text{analyse}$ (les deux mots partagent
le préfixe \og a \fg, puis \og l \fg{} précède \og n \fg{} dans l'alphabet).
\end{exemple}

%\begin{application}
%L'ordre lexicographique est utilisé pour trier des mots dans un dictionnaire,
%des noms dans un annuaire, ou des chaînes de caractères en informatique
%(fonction \texttt{sort} appliquée à des chaînes). Les relations d'ordre
%partiel modélisent, par exemple, les \textbf{dépendances entre tâches} dans
%la gestion de projet (une tâche \og précède \fg{} une autre) ou les relations
%d'\textbf{héritage} entre classes en programmation orientée objet.
%\end{application}
%
\section{Résumé du chapitre}

\begin{resume}
\begin{itemize}
  \item Un \textbf{ensemble} est une collection d'objets ; $x \in A$
  (appartenance) et $B \subset A$ (inclusion) sont deux notions distinctes.
  \item Opérations : $A\cap B$ (intersection), $A\cup B$ (union), $A\setminus
  B$ (différence), $C_E A$ (complémentaire) ; toutes se traduisent en
  connecteurs logiques ($\wedge,\vee,\overline{\ \cdot\ }$) sur les prédicats
  définissant $A$ et $B$.
  \item $\mathcal{P}(A)$ est l'ensemble des parties de $A$ (il a $2^n$
  éléments si $|A|=n$).
  \item Une \textbf{partition} découpe un ensemble en parties non vides,
  disjointes deux à deux, dont la réunion redonne l'ensemble entier.
  \item $A \times B = \{(a,b) : a\in A, b\in B\}$ : le produit cartésien
  n'est \textbf{pas commutatif} en général.
  \item Une \textbf{relation} sur $A$ peut être réflexive, symétrique,
  antisymétrique, transitive (ces deux dernières ne sont pas opposées).
  \item \textbf{Relation d'équivalence} = réflexive + symétrique +
  transitive $\Rightarrow$ les classes d'équivalence forment une
  \textbf{partition}.
  \item \textbf{Relation d'ordre} = réflexive + antisymétrique + transitive ;
  elle peut être totale (tous comparables) ou partielle (certains
  incomparables).
\end{itemize}
\end{resume}


\chapter{Applications}

\section{Introduction}


Ce chapitre a pour objectifs :
\begin{itemize}
  \item de définir rigoureusement une \textbf{application}, son
  \textbf{graphe}, l'\textbf{image directe} et l'\textbf{image réciproque}
  d'une partie ;
  \item de définir la \textbf{composition} de deux applications et d'en
  étudier les propriétés ;
  \item de caractériser les applications \textbf{injectives},
  \textbf{surjectives}, \textbf{bijectives}, et l'application
  \textbf{réciproque} d'une bijection ;
  \item d'établir le théorème de \textbf{décomposition canonique} d'une
  application, qui relie directement ce chapitre à la notion de relation
  d'équivalence vue au Chapitre 2.
\end{itemize}

\begin{application}
Les applications sont omniprésentes : en \textbf{informatique}, une fonction
de programmation est une application (à condition d'être bien définie pour
toute entrée valide) ; en \textbf{cryptographie}, le chiffrement doit être une
\textbf{bijection} pour être déchiffrable (application réciproque = clé de
déchiffrement) ; en \textbf{apprentissage automatique}, un modèle de
classification est une application de l'espace des données vers un ensemble
fini d'étiquettes, rarement injective.
\end{application}

\section{Prérequis}

Ce chapitre utilise abondamment :
\begin{itemize}
  \item le vocabulaire logique du Chapitre 1 (propositions, quantificateurs,
  contraposée) ;
  \item les notions d'ensemble, de produit cartésien, de relation binaire et,
  surtout, de \textbf{relation d'équivalence} et de \textbf{classe
  d'équivalence} du Chapitre 2, indispensables pour la section 3.6
  (décomposition canonique).
\end{itemize}

\section{Notion d'application}

\begin{definition}
On appelle \textbf{application} d'un ensemble $A$ dans un ensemble $B$ toute
correspondance $f$ qui associe, à \textbf{chaque} élément $x$ de $A$, \textbf{un
et un seul} élément $y$ de $B$. On dit que $A$ est l'\textbf{ensemble de
départ} (ou \textbf{source}) et $B$ l'\textbf{ensemble d'arrivée} (ou
\textbf{but}). L'élément $y$ associé à $x$ par $f$ s'appelle l'\textbf{image}
de $x$ par $f$, notée $f(x)$ (donc $y=f(x)$). On note :
\[
f : A \longrightarrow B, \qquad x \longmapsto y=f(x).
\]
Si $y=f(x)$, on dit que $x$ est un \textbf{antécédent} de $y$ par $f$.
\end{definition}

\begin{remarque}
Une image $y$ peut avoir \textbf{plusieurs} antécédents, mais un antécédent
$x$ ne peut jamais avoir \textbf{deux images distinctes} : c'est exactement
la condition \og un et un seul \fg{} de la définition.
\end{remarque}

\begin{exemple}[3.1]
\begin{enumerate}
  \item La correspondance $f$ qui associe à l'entier relatif $n$ l'entier
  naturel $n^2$ est une application de $\mathbb{Z}$ dans $\mathbb{N}$ :
  $f(n)=n^2$. On a $f(1)=1$, $f(-2)=4$.
  \item La correspondance $h$ qui associe au réel $x$ le réel
  $\dfrac{1}{x-1}$ \textbf{n'est pas} une application de $\mathbb{R}$ dans
  $\mathbb{R}$ : $1$ n'a pas d'image par $h$ (division par $0$). C'est en
  revanche une application de $\mathbb{R}\setminus\{1\}$ dans $\mathbb{R}$.
  \item La correspondance qui associe à chaque mois le nombre de jours
  possible de ce mois \textbf{n'est pas} une application de l'ensemble $M$
  des mois vers $\mathbb{N}$ : elle associe à \og février \fg{} les
  \emph{deux} valeurs $28$ et $29$, ce qui viole l'unicité de l'image.
  \item L'application \textbf{identique} de $A$, notée $\mathrm{Id}_A : A \to
  A$, associe à chaque élément $x$ l'élément $x$ lui-même :
  $\mathrm{Id}_A(x)=x$.
\end{enumerate}
\end{exemple}

\begin{definition}[3.2 -- Graphe, image directe, image réciproque]
Soit $f : A \to B$ une application.
\begin{enumerate}
  \item Le \textbf{graphe} de $f$ est $G_f = \{(x,f(x)) : x \in A\} \subset
  A\times B$.
  \item Si $A' \subset A$, l'\textbf{image directe} de $A'$ par $f$ est
  $f(A') = \{f(x) : x \in A'\}$. Le cas $A'=A$ donne
  $f(A)$, noté $\mathrm{Im}\,f$ et appelé \textbf{image} de $f$.
  \item Si $B' \subset B$, l'\textbf{image réciproque} de $B'$ par $f$ est
  $f^{-1}(B') = \{x \in A : f(x) \in B'\}$.
\end{enumerate}
\end{definition}


\begin{exemple}[3.2]
Soit $f : \mathbb{Z} \to \mathbb{N}$, $f(n) = n^2$.
\begin{enumerate}
  \item $f(\{-3,-1,0,3,5\}) = \{9,1,0,9,25\} = \{0,1,9,25\}$ (on ne répète
  pas les valeurs identiques dans un ensemble).
  \item $f^{-1}(\{4,5,6\}) = \{n \in \mathbb{Z} : n^2 \in \{4,5,6\}\}$.
  Cherchons les entiers dont le carré vaut $4$, $5$ ou $6$ : seul $n^2=4$ a
  des solutions entières, $n=2$ ou $n=-2$. Donc $f^{-1}(\{4,5,6\}) =
  \{-2,2\}$.
\end{enumerate}
\end{exemple}

\begin{proposition}[3.1 -- Propriétés des images directes et réciproques]
Soit $f : A \to B$ une application, $A_1,A_2 \subset A$ et $B_1,B_2\subset B$.
On a :
\begin{enumerate}[label=(\arabic*)]
  \item $f(A_1 \cup A_2) = f(A_1) \cup f(A_2)$ ;
  \item $f(A_1 \cap A_2) \subset f(A_1) \cap f(A_2)$ (inclusion, en général
  \textbf{stricte}) ;
  \item $f^{-1}(B_1 \cup B_2) = f^{-1}(B_1) \cup f^{-1}(B_2)$ ;
  \item $f^{-1}(B_1 \cap B_2) = f^{-1}(B_1) \cap f^{-1}(B_2)$.
\end{enumerate}
\end{proposition}

\begin{proof}[Démonstration de (2) et contre-exemple pour l'égalité]
Montrons $f(A_1\cap A_2) \subset f(A_1)\cap f(A_2)$. Soit $y \in f(A_1\cap
A_2)$ : il existe $x \in A_1 \cap A_2$ tel que $y=f(x)$. Comme $x\in A_1$,
$y=f(x) \in f(A_1)$ ; comme $x \in A_2$, $y=f(x)\in f(A_2)$. Donc $y \in
f(A_1)\cap f(A_2)$.

\textbf{L'inclusion réciproque est fausse en général.} Prenons $f:\mathbb{R}
\to \mathbb{R}$, $f(x)=x^2$, $A_1=\{-1\}$, $A_2=\{1\}$. Alors $A_1\cap
A_2=\varnothing$, donc $f(A_1\cap A_2)=\varnothing$. Mais $f(A_1)=\{1\}$ et
$f(A_2)=\{1\}$, donc $f(A_1)\cap f(A_2)=\{1\} \neq \varnothing$. L'inclusion
est donc bien stricte dans cet exemple : $f(A_1\cap A_2) \subsetneq
f(A_1)\cap f(A_2)$.
\end{proof}

\begin{remarque}[Pourquoi l'image réciproque se comporte mieux que l'image directe]
Les égalités (3) et (4) pour l'image réciproque, contrairement à
l'inclusion simple (2) pour l'image directe, s'expliquent ainsi : dire
\og $x \in f^{-1}(B_1\cap B_2)$ \fg{} équivaut exactement à
\og $f(x)\in B_1 \wedge f(x) \in B_2$ \fg, une conjonction de deux
conditions \emph{sur le même $x$}. Pour l'image directe, en revanche, un
élément de $f(A_1)\cap f(A_2)$ peut provenir d'un antécédent \emph{dans
$A_1$} et d'un \emph{autre} antécédent \emph{dans $A_2$} : c'est cette
possible \og dissociation des antécédents \fg{} qui casse l'égalité.
\end{remarque}

\section{Opérations sur les applications}

\subsection{Égalité de deux applications}

\begin{definition}[3.3]
Deux applications $f_1 : A_1 \to B_1$ et $f_2 : A_2 \to B_2$ sont
\textbf{égales} si $A_1=A_2$, $B_1=B_2$, et $\forall x \in A_1,\ f_1(x) =
f_2(x)$. On note alors $f_1=f_2$.
\end{definition}


\subsection{Composition des applications}


\begin{definition}[3.4]
Soient $f : A \to B$ et $g : B \to C$ deux applications. La \textbf{composée}
de $f$ suivie de $g$ est l'application notée $g\circ f$, définie par :
\[
g \circ f : A \longrightarrow C, \qquad \forall x \in A,\ (g\circ f)(x) = g(f(x)).
\]
\end{definition}


\begin{exemple}[3.3]
Soient $f : \mathbb{R} \to \mathbb{R}$, $f(x)=x+1$, et $g : \mathbb{R} \to
\mathbb{R}$, $g(x)=2x$. Calculons les deux composées :
\[
(g\circ f)(x) = g(f(x)) = g(x+1) = 2(x+1) = 2x+2,
\]
\[
(f\circ g)(x) = f(g(x)) = f(2x) = 2x+1.
\]
On constate que $g\circ f \neq f \circ g$ (par exemple, en $x=0$ :
$(g\circ f)(0)=2$ alors que $(f\circ g)(0)=1$) : \textbf{la composition
n'est pas commutative} en général.
\end{exemple}

\begin{proposition}[3.2]
Soient $f:A\to B$, $g:B\to C$, $h:C\to D$ trois applications. Alors :
\begin{enumerate}
  \item (\textbf{Associativité}) $(h\circ g)\circ f = h \circ (g\circ f)$.
  \item $f \circ \mathrm{Id}_A = f = \mathrm{Id}_B \circ f$.
\end{enumerate}
\end{proposition}

\begin{proof}
(1) Pour tout $x\in A$ :
\[
\big((h\circ g)\circ f\big)(x) = (h\circ g)(f(x)) = h\big(g(f(x))\big),
\]
\[
\big(h\circ (g\circ f)\big)(x) = h\big((g\circ f)(x)\big) = h\big(g(f(x))\big).
\]
Les deux expressions coïncident pour tout $x \in A$ : les deux applications
sont égales.

(2) Pour tout $x\in A$, $(f\circ \mathrm{Id}_A)(x) = f(\mathrm{Id}_A(x)) =
f(x)$, donc $f\circ\mathrm{Id}_A=f$. De même,
$(\mathrm{Id}_B\circ f)(x)=\mathrm{Id}_B(f(x))=f(x)$, donc
$\mathrm{Id}_B\circ f = f$.
\end{proof}

\section{Applications injectives, surjectives, bijectives}


\begin{definition}[3.5 -- Injectivité]
Une application $f : A \to B$ est \textbf{injective} si elle n'associe
jamais la même image à deux éléments \textbf{distincts} :
\[
f \text{ est injective} \iff \forall (x,x') \in A\times A,\ \big(f(x)=f(x')
\Rightarrow x=x'\big).
\]
\end{definition}

\begin{remarque}[Reformulation utile par contraposée]
D'après la loi de contraposée (Théorème 1.1), la définition équivaut à :
\[
f \text{ est injective} \iff \forall (x,x')\in A\times A,\ \big(x\neq x'
\Rightarrow f(x)\neq f(x')\big),
\]
c'est-à-dire : des éléments distincts ont toujours des images distinctes.
C'est souvent \emph{cette} formulation qu'on utilise pour \emph{montrer
qu'une application n'est pas injective} : il suffit d'exhiber $x\neq x'$
avec $f(x)=f(x')$.
\end{remarque}

\begin{definition}[3.6 -- Surjectivité]
Une application $f : A \to B$ est \textbf{surjective} si tout élément de
$B$ possède au moins un antécédent dans $A$ :
\[
f \text{ est surjective} \iff \forall y \in B,\ \exists x \in A,\ y=f(x)
\iff \mathrm{Im}\,f = B.
\]
\end{definition}

\begin{definition}[3.7 -- Bijectivité]
Une application $f : A \to B$ est \textbf{bijective} si elle est à la fois
injective et surjective.
\end{definition}

\begin{exemple}[3.4]
Soit $h : \mathbb{R} \to \mathbb{R}$, $h(x)=3x-1$.

\textbf{Injectivité :} soient $x,x' \in \mathbb{R}$ tels que $h(x)=h(x')$ :
\[
3x-1=3x'-1 \ \Longrightarrow\ 3x=3x' \ \Longrightarrow\ x=x'.
\]
Donc $h$ est injective.

\textbf{Surjectivité :} soit $y\in\mathbb{R}$. Cherchons $x\in\mathbb{R}$ tel
que $y=h(x)$ :
\[
y=3x-1 \iff y+1=3x \iff x=\frac{y+1}{3}.
\]
Pour tout $y\in\mathbb{R}$, $x=\dfrac{y+1}{3}$ est bien un réel et vérifie
$h(x)=y$. Donc $h$ est surjective.

$h$ est donc \textbf{bijective}.
\end{exemple}

\begin{exemple}[3.5]
Soit $g : \mathbb{Z} \to \mathbb{N}$, $g(n)=n^2$.

\textbf{$g$ n'est pas injective} : $g(2)=4=g(-2)$, mais $2\neq -2$.

\textbf{$g$ n'est pas surjective} : cherchons un antécédent à $5$. Il
faudrait $n^2=5$ avec $n\in\mathbb{Z}$, ce qui est impossible ($2^2=4<5<9=3^2$,
aucun entier n'a pour carré $5$). Donc $5 \notin \mathrm{Im}\,g$ : $g$ n'est
pas surjective.
\end{exemple}

\begin{proposition}[3.3 -- Composition et injectivité/surjectivité]
Soient $f : A \to B$ et $g : B \to C$ deux applications.
\begin{enumerate}
  \item Si $f$ et $g$ sont injectives, alors $g\circ f$ est injective.
  \item Si $f$ et $g$ sont surjectives, alors $g\circ f$ est surjective.
  \item Si $g\circ f$ est injective, alors $f$ est injective.
  \item Si $g\circ f$ est surjective, alors $g$ est surjective.
\end{enumerate}
\end{proposition}

\begin{proof}
(1) Supposons $f$ et $g$ injectives. Soient $x,x'\in A$ tels que $(g\circ
f)(x)=(g\circ f)(x')$, c'est-à-dire $g(f(x))=g(f(x'))$. Comme $g$ est
injective, on en déduit $f(x)=f(x')$. Comme $f$ est injective, on en déduit
$x=x'$. Donc $g\circ f$ est injective.

(2) Supposons $f$ et $g$ surjectives. Soit $z\in C$. Comme $g$ est
surjective, il existe $y\in B$ tel que $z=g(y)$. Comme $f$ est surjective,
il existe $x\in A$ tel que $y=f(x)$. Alors
$(g\circ f)(x) = g(f(x)) = g(y) = z$. Donc $g\circ f$ est surjective.

(3) Supposons $g\circ f$ injective. Soient $x,x'\in A$ tels que $f(x)=f(x')$.
Alors $g(f(x))=g(f(x'))$, c'est-à-dire $(g\circ f)(x)=(g\circ f)(x')$. Comme
$g\circ f$ est injective, on en déduit $x=x'$. Donc $f$ est injective.

(4) Supposons $g\circ f$ surjective. Soit $z\in C$. Comme $g\circ f$ est
surjective, il existe $x\in A$ tel que $z=(g\circ f)(x)=g(f(x))$. En posant
$y=f(x)\in B$, on a $z=g(y)$ : $z$ possède bien un antécédent par $g$. Donc
$g$ est surjective.
\end{proof}


\section{Application réciproque d'une bijection}


\begin{definition}[3.8]
Soit $f : A \to B$ une application bijective. On appelle \textbf{application
réciproque} de $f$ l'application, notée $f^{-1}$, définie par
\[
f^{-1} : B \longrightarrow A, \qquad f^{-1}(y) = x \quad \text{où } x
\text{ est l'unique antécédent de } y \text{ par } f \text{ (i.e. } f(x)=y).
\]
\end{definition}

\begin{theoreme}[3.1]
Si $f : A \to B$ est bijective, alors $f^{-1}$ est également bijective,
$(f^{-1})^{-1} = f$, et
\[
f \circ f^{-1} = \mathrm{Id}_B, \qquad f^{-1}\circ f = \mathrm{Id}_A.
\]
\end{theoreme}

\begin{proof}
\textbf{Les deux égalités de composition.} Soit $y\in B$. Par définition,
$f^{-1}(y)$ est l'unique $x\in A$ tel que $f(x)=y$. Donc
$(f\circ f^{-1})(y) = f(f^{-1}(y)) = f(x) = y$. Ceci vaut pour tout $y\in B$,
donc $f\circ f^{-1} = \mathrm{Id}_B$.

Soit $x\in A$, posons $y=f(x)$. Par définition de $f^{-1}$, $f^{-1}(y)$ est
l'unique antécédent de $y$, or $x$ en est justement un (puisque $f(x)=y$) ;
par unicité, $f^{-1}(y)=x$. Donc $(f^{-1}\circ f)(x) = f^{-1}(f(x)) =
f^{-1}(y) = x$. Ceci vaut pour tout $x\in A$, donc $f^{-1}\circ f =
\mathrm{Id}_A$.

\textbf{Bijectivité de $f^{-1}$.} \emph{Injectivité :} soient $y,y'\in B$
tels que $f^{-1}(y)=f^{-1}(y')$. En appliquant $f$ aux deux membres :
$f(f^{-1}(y))=f(f^{-1}(y'))$, c'est-à-dire, d'après l'égalité établie
ci-dessus, $y=y'$. Donc $f^{-1}$ est injective.

\emph{Surjectivité :} soit $x\in A$. Posons $y=f(x)\in B$. D'après ce qui
précède, $f^{-1}(y)=x$ : $x$ possède bien un antécédent (à savoir $y$) par
$f^{-1}$. Donc $f^{-1}$ est surjective.

$f^{-1}$ est donc bijective, et $(f^{-1})^{-1}=f$ découle immédiatement de la
définition (l'unique antécédent de $x$ par $f^{-1}$ est précisément
$f(x)$).
\end{proof}

\begin{exemple}[3.6]
Reprenons $h : \mathbb{R} \to \mathbb{R}$, $h(x)=3x-1$, bijective (Exemple
3.4). D'après le calcul effectué (résolution de $y=h(x)$), son application
réciproque est
\[
h^{-1} : \mathbb{R}\to\mathbb{R}, \qquad h^{-1}(y) = \frac{y+1}{3}.
\]
\end{exemple}

\begin{proposition}[3.4]
Si $f:A\to B$ et $g:B\to C$ sont deux bijections, alors $g\circ f$ est une
bijection et $(g\circ f)^{-1} = f^{-1}\circ g^{-1}$.
\end{proposition}

\begin{proof}
La bijectivité de $g\circ f$ découle immédiatement de la Proposition 3.3
(points 1 et 2, $f,g$ étant à la fois injectives et surjectives). Pour la
formule de l'inverse, calculons $(g\circ f)\circ(f^{-1}\circ g^{-1})$ :
\[
(g\circ f)\circ(f^{-1}\circ g^{-1}) = g\circ (f\circ f^{-1}) \circ g^{-1}
= g \circ \mathrm{Id}_B \circ g^{-1} = g\circ g^{-1} = \mathrm{Id}_C,
\]
en utilisant l'associativité de la composition (Proposition 3.2) et le
Théorème 3.1. De même, on vérifie $(f^{-1}\circ g^{-1})\circ(g\circ f) =
\mathrm{Id}_A$. Ces deux égalités caractérisent $f^{-1}\circ g^{-1}$ comme
étant la réciproque de $g\circ f$.
\end{proof}

\section{Décomposition canonique d'une application}


\begin{definition}[3.9 -- Injection canonique]
Si $A' \subset A$, l'\textbf{injection canonique} de $A'$ dans $A$ est
l'application $i : A' \to A$, $i(x)=x$ (elle est bien injective, puisque
$i(x)=i(x') \Rightarrow x=x'$ trivialement).
\end{definition}

\begin{definition}[3.10 -- Relation associée à une application]
Soit $f : E \to F$ une application. On définit la relation $\mathcal{R}_f$
sur $E$ par :
\[
\forall x,y \in E,\quad x\,\mathcal{R}_f\,y \iff f(x)=f(y).
\]
\end{definition}

\begin{proposition}[3.5]
$\mathcal{R}_f$ est une relation d'équivalence sur $E$.
\end{proposition}

\begin{proof}
Il suffit de traduire les propriétés de l'égalité dans $F$ (voir Exercice
E4 du Chapitre 2).

\textbf{Réflexivité :} pour tout $x\in E$, $f(x)=f(x)$, donc
$x\,\mathcal{R}_f\,x$.

\textbf{Symétrie :} si $f(x)=f(y)$, alors $f(y)=f(x)$ (symétrie de
l'égalité), donc $y\,\mathcal{R}_f\,x$.

\textbf{Transitivité :} si $f(x)=f(y)$ et $f(y)=f(z)$, alors $f(x)=f(z)$
(transitivité de l'égalité), donc $x\,\mathcal{R}_f\,z$.
\end{proof}

D'après le Théorème 2.1 (Chapitre 2), les classes d'équivalence de
$\mathcal{R}_f$ forment donc une partition de $E$. On peut noter $E/f$
l'ensemble quotient $E/\mathcal{R}_f$.

\begin{theoreme}[3.2 -- Théorème de décomposition canonique]
Soit $f : E \to F$ une application. Il existe :
\begin{itemize}
  \item une application \textbf{surjective} $s : E \to E/\mathcal{R}_f$,
  dite \textbf{surjection canonique}, définie par $s(x) = \dot x$ ;
  \item une application \textbf{bijective} $\bar f : E/\mathcal{R}_f \to
  \mathrm{Im}\,f$, définie par $\bar f(\dot x) = f(x)$ ;
  \item l'\textbf{injection canonique} $i : \mathrm{Im}\,f \to F$,
\end{itemize}
telles que $f = i \circ \bar f \circ s$. Autrement dit, le diagramme suivant
\og commute \fg{} :
\[
E \xrightarrow{\ f\ } F \qquad = \qquad
E \xrightarrow{\ s\ } E/\mathcal{R}_f \xrightarrow{\ \bar f\ } \mathrm{Im}\,f
\xrightarrow{\ i\ } F.
\]
\end{theoreme}

\begin{proof}
\textbf{$s$ est bien définie et surjective.} $s(x)=\dot x$ associe à $x$ sa
classe d'équivalence, qui existe et est unique pour tout $x\in E$ : $s$ est
une application. Elle est surjective par construction : toute classe
$\dot x \in E/\mathcal{R}_f$ est bien l'image de $x$ par $s$.

\textbf{$\bar f$ est bien définie (ne dépend pas du représentant choisi).}
Il faut vérifier que si $\dot x = \dot y$ (même classe), alors $f(x)=f(y)$
(même image), sans quoi $\bar f(\dot x)$ ne serait pas définie sans
ambiguïté. Or $\dot x = \dot y$ signifie exactement $x\,\mathcal{R}_f\,y$,
c'est-à-dire, par définition même de $\mathcal{R}_f$, $f(x)=f(y)$. La
définition de $\bar f$ est donc cohérente (\og bien posée \fg).

\textbf{$\bar f$ est injective.} Soient $\dot x, \dot y \in E/\mathcal{R}_f$
tels que $\bar f(\dot x) = \bar f(\dot y)$, c'est-à-dire $f(x)=f(y)$. Par
définition de $\mathcal{R}_f$, cela signifie $x\,\mathcal{R}_f\,y$,
c'est-à-dire $\dot x = \dot y$. Donc $\bar f$ est injective.

\textbf{$\bar f$ est surjective (sur $\mathrm{Im}\,f$).} Soit $z \in
\mathrm{Im}\,f$. Par définition de l'image, il existe $x\in E$ tel que
$z=f(x)$. Alors $\bar f(\dot x) = f(x) = z$ : $z$ possède bien un
antécédent par $\bar f$. Donc $\bar f$ est surjective sur $\mathrm{Im}\,f$.

$\bar f$ est donc bijective.

\textbf{Vérification de $f = i\circ \bar f \circ s$.} Soit $x\in E$. On
calcule :
\[
(i\circ \bar f \circ s)(x) = i\big(\bar f(s(x))\big) = i\big(\bar f(\dot
x)\big) = i\big(f(x)\big) = f(x),
\]
puisque $i$ est l'injection canonique (elle laisse les éléments inchangés).
Ceci vaut pour tout $x \in E$ : les deux applications $f$ et $i\circ \bar f
\circ s$ sont égales.
\end{proof}

\begin{remarque}[Interprétation]
Ce théorème dit, en substance : \textbf{toute application se décompose en
une surjection suivie d'une bijection suivie d'une injection}. C'est un
résultat fondamental de \og classification \fg : $\bar f$ réalise, en
quelque sorte, la \og vraie \fg{} correspondance bijective cachée derrière
$f$, une fois qu'on a \og recollé \fg{} entre eux les éléments ayant la même
image.
\end{remarque}

\begin{exemple}[3.7]
Soit $f : \mathbb{R} \to \mathbb{R}$, $f(x)=x^2$.

La relation $\mathcal{R}_f$ associée est $x\,\mathcal{R}_f\,y \iff x^2=y^2
\iff (x=y) \vee (x=-y)$ (voir aussi l'exemple analogue au Chapitre 1,
section 1.2). Les classes d'équivalence sont $\dot 0 = \{0\}$ et, pour
$a\neq0$, $\dot a = \{a,-a\}$.

On peut identifier l'ensemble quotient $\mathbb{R}/\mathcal{R}_f$ à
$\mathbb{R}^+$ en choisissant, dans chaque classe $\{a,-a\}$ (ou $\{0\}$),
le représentant positif $|a|$. L'application $\bar f$ devient alors
$\bar f : \mathbb{R}^+ \to \mathrm{Im}\,f = \mathbb{R}^+$, $\bar f(t) = t^2$,
qui est bien \textbf{bijective} (c'est la fonction carré restreinte aux
réels positifs, strictement croissante donc injective, et clairement
surjective sur $\mathbb{R}^+$). Enfin, $i : \mathbb{R}^+ \to \mathbb{R}$ est
l'injection canonique.

On retrouve ainsi, de façon très concrète, pourquoi $f(x)=x^2$
\og devient bijective \fg{} dès qu'on la restreint aux réels positifs : c'est
exactement le contenu du théorème de décomposition canonique.
\end{exemple}


\chapter{Notions de groupes, d'anneaux et de corps}

\section{Introduction}


Ce chapitre a pour objectifs :
\begin{itemize}
  \item de définir une \textbf{loi de composition interne} et ses
  propriétés remarquables (associativité, élément neutre, éléments
  inversibles, commutativité) ;
  \item de définir la structure de \textbf{groupe}, et d'introduire la
  notion de \textbf{sous-groupe} ;
  \item d'étudier les \textbf{homomorphismes de groupes}, qui \og
  transportent \fg{} la structure d'un groupe vers un autre ;
  \item de définir les structures d'\textbf{anneau} et de \textbf{corps},
  qui combinent deux lois de composition interne.
\end{itemize}

\begin{application}
Les structures algébriques sont au c\oe ur de la \textbf{cryptographie
moderne} (groupes de points sur une courbe elliptique, corps finis
$\mathbb{Z}/p\mathbb{Z}$ utilisés dans RSA), de la \textbf{théorie des codes
correcteurs d'erreurs} (transmission de données), de la \textbf{physique
théorique} (les groupes de symétrie décrivent les lois de conservation), et
de l'\textbf{informatique graphique} (groupes de rotations et de
symétries).
\end{application}

\section{Prérequis}

Ce chapitre s'appuie sur le vocabulaire logique (Chapitre 1), la notion
d'ensemble et de relation d'équivalence (Chapitre 2 -- en particulier les
classes modulo $n$, qui reviendront en exemple), et sur la notion
d'application, injective/surjective/bijective (Chapitre 3), indispensable
pour définir les homomorphismes.

\section{Loi de composition interne}


\begin{definition}[4.1]
On appelle \textbf{loi de composition interne} (LCI), ou \textbf{opération
binaire}, sur un ensemble non vide $E$, toute application $*$ de $E\times E$
dans $E$. L'image $*(x,y)$ est notée $x*y$. Autrement dit :
\[
\big(* \text{ est une LCI sur } E\big) \iff
\begin{cases}
\forall x,y \in E,\ x*y \in E & \text{(stabilité)}\\
\forall x,y,x',y' \in E,\ (x=x' \wedge y=y') \Rightarrow x*y=x'*y' &
\text{(compatibilité)}
\end{cases}
\]
\end{definition}

\begin{remarque}
La deuxième condition (\textbf{compatibilité}) est souvent oubliée, car elle
paraît évidente : elle garantit simplement que le résultat $x*y$ ne dépend
que \emph{des valeurs} de $x$ et $y$, pas de la façon dont on les a
désignés. C'est précisément cette propriété qui permettra de définir des
opérations sur un \textbf{ensemble quotient} (Chapitre 2) : par exemple,
l'addition $\dot a + \dot b := \dot{(a+b)}$ sur $\mathbb{Z}/n\mathbb{Z}$
n'est légitime que parce qu'elle ne dépend pas du représentant $a$ (ou $b$)
choisi dans chaque classe -- nous y reviendrons en exemple.
\end{remarque}

\begin{exemple}[4.1]
\begin{enumerate}
  \item L'addition $+$ et la multiplication $\cdot$ usuelles sont des LCI
  sur $\mathbb{N}$, $\mathbb{Z}$, $\mathbb{Q}$, $\mathbb{R}$, $\mathbb{C}$.
  \item La soustraction $-$ est une LCI sur $\mathbb{Z}$, $\mathbb{Q}$,
  $\mathbb{R}$, $\mathbb{C}$, mais \textbf{pas} sur $\mathbb{N}$ : par
  exemple $2,5\in\mathbb{N}$ mais $2-5=-3\notin\mathbb{N}$ (défaut de
  stabilité).
  \item La composition $\circ$ est une LCI sur l'ensemble
  $\mathcal{A}(E,E)$ des applications de $E$ dans $E$ (Chapitre 3).
  \item L'intersection $\cap$ est une LCI sur $\mathcal{P}(E)$, l'ensemble
  des parties de $E$ (Chapitre 2).
\end{enumerate}
\end{exemple}

\begin{exemple}[4.2 -- Une opération bien définie sur un ensemble quotient]
Reprenons la relation de congruence modulo $n$ sur $\mathbb{Z}$ (Chapitre 2,
Exemple 2.8) et son ensemble quotient $\mathbb{Z}/n\mathbb{Z}$. On souhaite
définir une addition $\dot a \oplus \dot b := \dot{(a+b)}$. Pour que ce soit
une LCI \textbf{bien définie}, il faut vérifier la compatibilité : si
$\dot a = \dot{a'}$ et $\dot b = \dot{b'}$ (mêmes classes, mais
représentants différents), alors $\dot{(a+b)}$ doit être égale à
$\dot{(a'+b')}$.

Or $\dot a=\dot{a'}$ signifie $n$ divise $a-a'$, et $\dot b=\dot{b'}$
signifie $n$ divise $b-b'$. En additionnant, $n$ divise
$(a-a')+(b-b') = (a+b)-(a'+b')$, donc $\dot{(a+b)} = \dot{(a'+b')}$. La
définition est donc cohérente : $\oplus$ est bien une LCI sur
$\mathbb{Z}/n\mathbb{Z}$.
\end{exemple}

\section{Propriétés d'une loi de composition interne}


\begin{definition}[4.2]
Soit $*$ une LCI sur un ensemble non vide $E$.
\begin{enumerate}
  \item $*$ est \textbf{associative} si
  $\forall x,y,z\in E,\ (x*y)*z = x*(y*z)$.
  \item $*$ admet un \textbf{élément neutre} s'il existe $e\in E$ tel que
  $\forall x\in E,\ (x*e=x) \wedge (e*x=x)$.
  \item Si $*$ admet un élément neutre $e$, un élément $x\in E$ est dit
  \textbf{inversible} (ou \textbf{symétrisable}) s'il existe $x'\in E$ tel
  que $(x*x'=e)\wedge(x'*x=e)$. L'élément $x'$ est appelé
  \textbf{inverse} (ou \textbf{symétrique}) de $x$, noté $x^{-1}$.
  \item $*$ est \textbf{commutative} si $\forall x,y\in E,\ x*y=y*x$.
\end{enumerate}
\end{definition}

\begin{remarque}[Pourquoi la parenthèse compte]
Sans associativité, une expression comme $x*y*z$ n'aurait \textbf{aucun
sens} : il faudrait toujours préciser si l'on calcule $(x*y)*z$ ou
$x*(y*z)$, qui pourraient différer. L'associativité est ce qui permet
d'omettre les parenthèses en toute sécurité.
\end{remarque}

\begin{theoreme}[4.1 -- Unicité de l'élément neutre]
Si une LCI $*$ sur $E$ admet un élément neutre, celui-ci est \textbf{unique}.
\end{theoreme}

\begin{proof}
Supposons que $e$ et $e'$ soient deux éléments neutres de $*$. Comme $e$
est neutre, en particulier appliqué à $e'$ : $e*e' = e'$. Comme $e'$ est
neutre, en particulier appliqué à $e$ : $e*e' = e$. En comparant les deux
égalités, $e=e*e'=e'$, donc $e=e'$ : l'élément neutre, s'il existe, est
unique.
\end{proof}

\begin{theoreme}[4.2 -- Unicité de l'inverse (loi associative)]
Soit $*$ une LCI \textbf{associative} sur $E$, admettant un élément neutre
$e$. Si $x\in E$ est inversible, son inverse est \textbf{unique}.
\end{theoreme}

\begin{proof}
Supposons que $x'$ et $x''$ soient deux inverses de $x$ : $x*x'=x'*x=e$ et
$x*x''=x''*x=e$. Calculons $x'*(x*x'')$ de deux façons, en utilisant
l'associativité :
\[
x' = x'*e = x'*(x*x'') = (x'*x)*x'' = e*x'' = x''.
\]
Donc $x'=x''$ : l'inverse, s'il existe, est unique (\textbf{à condition que
$*$ soit associative}, hypothèse utilisée de façon essentielle au moment
d'appliquer l'associativité).
\end{proof}


\begin{exemple}[4.3]
Étudions la loi $*$ définie sur $\mathbb{R}$ par $x*y = x+y-1$.

\textbf{Associativité :} $(x*y)*z = (x+y-1)*z = (x+y-1)+z-1 = x+y+z-2$.
$x*(y*z) = x*(y+z-1) = x+(y+z-1)-1 = x+y+z-2$. Les deux expressions
coïncident : $*$ est associative.

\textbf{Élément neutre :} on cherche $e$ tel que $x*e=x$ pour tout $x$ :
$x+e-1=x \Rightarrow e=1$. On vérifie que $e=1$ convient aussi à gauche
(la loi étant commutative, voir ci-dessous) : $1*x=1+x-1=x$. L'élément
neutre est $e=1$.

\textbf{Inversibilité :} pour $x\in\mathbb{R}$, on cherche $x'$ tel que
$x*x'=1$ : $x+x'-1=1 \Rightarrow x'=2-x$. Tout réel $x$ admet donc un
inverse, $x^{-1}=2-x$.

\textbf{Commutativité :} $x*y=x+y-1=y+x-1=y*x$ : $*$ est commutative.
\end{exemple}

\section{Structure de groupe}


\begin{definition}[4.3]
Soit $*$ une LCI sur un ensemble non vide $G$. On dit que $(G,*)$ est un
\textbf{groupe} si :
\begin{enumerate}
  \item $*$ est \textbf{associative} ;
  \item $*$ admet un \textbf{élément neutre} $e$ ;
  \item \textbf{tout} élément de $G$ est \textbf{inversible} par rapport à
  $*$.
\end{enumerate}
Si de plus $*$ est \textbf{commutative}, on dit que le groupe est
\textbf{commutatif} (ou \textbf{abélien}).
\end{definition}

\begin{exemple}[4.4]
\begin{enumerate}
  \item $(\mathbb{Z},+)$, $(\mathbb{Q},+)$, $(\mathbb{R},+)$,
  $(\mathbb{C},+)$ sont des groupes commutatifs (élément neutre $0$,
  inverse de $x$ : $-x$).
  \item $(\mathbb{Q}^*,\cdot)$, $(\mathbb{R}^*,\cdot)$,
  $(\mathbb{C}^*,\cdot)$ sont des groupes commutatifs (élément neutre $1$,
  inverse de $x$ : $\dfrac{1}{x}$). On note $\mathbb{K}^*=\mathbb{K}
  \setminus\{0\}$.
  \item L'ensemble $\mathbb{R}$ muni de la loi $*$ de l'Exemple 4.3
  ($x*y=x+y-1$) est un groupe commutatif, d'élément neutre $1$ et
  d'inverse $x^{-1}=2-x$.
  \item $(\mathbb{N},+)$ \textbf{n'est pas} un groupe : $1\in\mathbb{N}$
  n'a pas d'inverse pour $+$ dans $\mathbb{N}$ (il faudrait $-1$, qui
  n'appartient pas à $\mathbb{N}$).
  \item $\big(\mathbb{Z}/n\mathbb{Z}, \oplus\big)$, avec l'addition
  $\oplus$ bien définie de l'Exemple 4.2, est un groupe commutatif
  d'élément neutre $\dot 0$ : c'est une conséquence directe des propriétés
  de l'addition dans $\mathbb{Z}$, transportées au quotient.
\end{enumerate}
\end{exemple}


\section{Sous-groupes}


\begin{definition}[4.4]
Soit $(G,*)$ un groupe et $H$ une partie de $G$. On dit que $(H,*)$ est un
\textbf{sous-groupe} de $(G,*)$ si $(H,*)$ est lui-même un groupe pour la
loi $*$ restreinte à $H$.
\end{definition}

\begin{theoreme}[4.3 -- Caractérisation pratique des sous-groupes]
Soit $(G,*)$ un groupe d'élément neutre $e$, et $H$ une partie de $G$.
Alors :
\[
\big((H,*) \text{ est un sous-groupe de } (G,*)\big) \iff
\begin{cases}
e \in H \\
\forall x,y \in H,\ x*y^{-1} \in H.
\end{cases}
\]
\end{theoreme}

\begin{proof}
\textbf{Sens direct.} Supposons $(H,*)$ sous-groupe de $(G,*)$. Comme
$(H,*)$ est un groupe, il admet un élément neutre $e_H$, et l'on montre
d'abord que $e_H=e$ : pour $x\in H\subset G$, $x*e_H=x$ (neutralité dans
$H$) et $x*e=x$ (neutralité dans $G$) ; en composant à gauche par
$x^{-1}$ (l'inverse de $x$ dans $G$) : $e_H = x^{-1}*x*e_H = x^{-1}*x =
e$. Donc $e\in H$.

Soient $x,y\in H$. Comme $(H,*)$ est un groupe, $y$ admet un inverse dans
$H$, qui coïncide avec son inverse $y^{-1}$ dans $G$ (par unicité de
l'inverse, Théorème 4.2), donc $y^{-1}\in H$. Comme $*$ est une LCI sur
$H$ (stabilité), $x*y^{-1}\in H$.

\textbf{Sens réciproque.} Supposons $e\in H$ et $\forall x,y\in H,\
x*y^{-1}\in H$. Montrons que $(H,*)$ est un groupe.

\emph{$H$ est non vide} (il contient $e$).

\emph{Tout élément de $H$ est inversible dans $H$.} Soit $y\in H$.
Appliquons l'hypothèse à $x=e\in H$ et $y\in H$ : $e*y^{-1} = y^{-1} \in
H$. Donc l'inverse (dans $G$) de tout élément de $H$ appartient encore à
$H$.

\emph{$*$ est une LCI sur $H$ (stabilité).} Soient $x,y\in H$. D'après le
point précédent, $y^{-1}\in H$. Appliquons l'hypothèse à $x\in H$ et
$y^{-1}\in H$ : $x*(y^{-1})^{-1} = x*y \in H$ (on utilise
$(y^{-1})^{-1}=y$, propriété de l'inverse dans un groupe). Donc $H$ est
stable par $*$.

\emph{$e$ est neutre dans $H$, $*$ est associative dans $H$.} Ces deux
propriétés sont héritées directement de $G$ (elles sont vraies pour
\emph{tous} les éléments de $G$, donc en particulier pour ceux de $H$).

Toutes les conditions d'un groupe sont donc vérifiées pour $(H,*)$ : c'est
un sous-groupe de $(G,*)$.
\end{proof}

\begin{exemple}[4.5]
\begin{enumerate}
  \item $\mathbb{Z}$ est une partie de $\mathbb{Q}$, et $(\mathbb{Q},+)$
  est un groupe. On a $0\in\mathbb{Z}$, et pour tous $x,y\in\mathbb{Z}$,
  $x-y=x+(-y)\in\mathbb{Z}$. Donc $(\mathbb{Z},+)$ est un sous-groupe de
  $(\mathbb{Q},+)$.
  \item Si $(G,*)$ est un groupe d'élément neutre $e$, alors $(\{e\},*)$
  et $(G,*)$ sont toujours des sous-groupes de $(G,*)$ : on les appelle
  les \textbf{sous-groupes triviaux}.
  \item $\mathbb{R}^{*+} = \{x\in\mathbb{R} : x>0\}$ est un sous-groupe de
  $(\mathbb{R}^*,\cdot)$ : $1\in\mathbb{R}^{*+}$, et pour $x,y\in
  \mathbb{R}^{*+}$, $x\cdot y^{-1} = \dfrac{x}{y} > 0$ (quotient de deux
  réels strictement positifs).
\end{enumerate}
\end{exemple}

\section{Homomorphismes de groupes}


\begin{definition}[4.5]
Soient $(G,*)$ et $(G',*')$ deux groupes. On appelle \textbf{homomorphisme}
de $(G,*)$ dans $(G',*')$ toute application $f : G \to G'$ telle que
\[
\forall x,y \in G,\quad f(x*y) = f(x) *' f(y).
\]
\end{definition}

\begin{exemple}[4.6]
L'application $h : \mathbb{R} \to \mathbb{R}^{*+}$, $h(x)=e^x$, est un
homomorphisme du groupe $(\mathbb{R},+)$ dans le groupe
$(\mathbb{R}^{*+},\cdot)$. En effet, pour $x,y\in\mathbb{R}$ :
\[
h(x+y) = e^{x+y} = e^x \cdot e^y = h(x)\cdot h(y).
\]
De même, l'application module $f:\mathbb{C}^*\to\mathbb{R}^*$,
$f(z)=|z|$, est un homomorphisme de $(\mathbb{C}^*,\cdot)$ dans
$(\mathbb{R}^*,\cdot)$, car $|z\cdot z'| = |z|\cdot|z'|$ pour tous
$z,z'\in\mathbb{C}^*$.
\end{exemple}

\begin{theoreme}[4.4]
Soit $f : G \to G'$ un homomorphisme du groupe $(G,*)$ dans le groupe
$(G',*')$, d'éléments neutres respectifs $e$ et $e'$. Alors :
\begin{enumerate}
  \item $f(e) = e'$.
  \item Pour tout $x\in G$, $f(x^{-1}) = \big(f(x)\big)^{-1}$.
\end{enumerate}
\end{theoreme}

\begin{proof}
(1) Calculons $f(e)*'f(e)$ de deux façons. D'une part, comme $f$ est un
homomorphisme, $f(e)*'f(e) = f(e*e) = f(e)$ (car $e*e=e$, $e$ étant
neutre). D'autre part, notons $f(e)^{-1}$ l'inverse de $f(e)$ dans $G'$
(qui existe puisque $(G',*')$ est un groupe). En composant l'égalité
$f(e)*'f(e)=f(e)$ à droite par $f(e)^{-1}$ :
\[
\big(f(e)*'f(e)\big) *' f(e)^{-1} = f(e) *' f(e)^{-1}
\ \Longrightarrow\
f(e) *' \big(f(e) *' f(e)^{-1}\big) = e'
\ \Longrightarrow\
f(e) *' e' = e'
\ \Longrightarrow\
f(e) = e'.
\]

(2) Soit $x\in G$. Comme $f$ est un homomorphisme :
\[
f(x^{-1}) *' f(x) = f(x^{-1}*x) = f(e) = e' \quad\text{(d'après (1))},
\]
\[
f(x) *' f(x^{-1}) = f(x*x^{-1}) = f(e) = e'.
\]
Ces deux égalités montrent exactement que $f(x^{-1})$ est l'inverse de
$f(x)$ dans $G'$ ; par unicité de l'inverse (Théorème 4.2),
$f(x^{-1}) = \big(f(x)\big)^{-1}$.
\end{proof}

\section{Structure d'anneau}


\begin{definition}[4.6]
Soit $A$ un ensemble non vide muni de deux LCI $*_1$ et $*_2$. On dit que
$(A,*_1,*_2)$ est un \textbf{anneau} si :
\begin{enumerate}
  \item $(A,*_1)$ est un \textbf{groupe commutatif} ;
  \item $*_2$ est \textbf{associative} ;
  \item $*_2$ est \textbf{distributive} par rapport à $*_1$ :
  $\forall x,y,z\in A$,
  \[
  x*_2(y*_1z) = (x*_2y)*_1(x*_2z) \qquad \text{et} \qquad
  (y*_1z)*_2x = (y*_2x)*_1(z*_2x).
  \]
\end{enumerate}
Si $*_2$ admet un élément neutre, on l'appelle \textbf{unité} et on dit que
l'anneau est \textbf{unitaire}. Si $*_2$ est commutative, l'anneau est dit
\textbf{commutatif}.
\end{definition}

\begin{remarque}
On note en général les deux lois d'un anneau $+_A$ et $\cdot_A$ (au lieu de
$*_1,*_2$), l'élément neutre de $+_A$ par $0_A$, et (s'il existe)
l'élément neutre de $\cdot_A$ par $1_A$.
\end{remarque}

\begin{exemple}[4.7]
$(\mathbb{Z},+,\cdot)$, $(\mathbb{Q},+,\cdot)$, $(\mathbb{R},+,\cdot)$,
$(\mathbb{C},+,\cdot)$ sont tous des anneaux commutatifs et unitaires
(unité $1$). Nous verrons au Chapitre 6 que $(K[X],+,\cdot)$, l'ensemble
des polynômes à coefficients dans un corps $K$, est également un anneau
commutatif et unitaire.
\end{exemple}

\begin{theoreme}[4.5 -- Règles de calcul dans un anneau]
Soit $(A,+_A,\cdot_A)$ un anneau d'élément neutre $0_A$ pour $+_A$. Alors :
\begin{enumerate}
  \item $\forall x\in A,\ x\cdot_A 0_A = 0_A = 0_A \cdot_A x$.
  \item $\forall x,y\in A,\ (-x)\cdot_A y = -(x\cdot_A y) = x\cdot_A(-y)$.
  \item $\forall x,y\in A,\ (-x)\cdot_A(-y) = x\cdot_A y$.
\end{enumerate}
\end{theoreme}

\begin{proof}
(1) Soit $x\in A$. On a, en utilisant la distributivité et le fait que
$0_A$ est neutre pour $+_A$ (donc $0_A+_A0_A=0_A$) :
\[
x\cdot_A 0_A = x\cdot_A(0_A+_A0_A) = (x\cdot_A0_A) +_A (x\cdot_A0_A).
\]
Notons $t = x\cdot_A 0_A$. On vient d'établir $t = t+_At$. En composant
avec $-t$ (l'inverse de $t$ pour $+_A$, qui existe car $(A,+_A)$ est un
groupe) : $t +_A(-t) = (t+_At)+_A(-t)$, soit $0_A = t +_A (t+_A(-t)) =
t+_A0_A = t$. Donc $t=0_A$, c'est-à-dire $x\cdot_A0_A=0_A$. On montre de
façon symétrique $0_A\cdot_Ax=0_A$.

(2) Soient $x,y\in A$. On calcule, en utilisant la distributivité :
\[
x\cdot_Ay +_A (-x)\cdot_Ay = \big(x+_A(-x)\big)\cdot_Ay = 0_A\cdot_Ay =
0_A \quad\text{(d'après (1))}.
\]
Cette égalité montre exactement que $(-x)\cdot_Ay$ est l'opposé (l'inverse
pour $+_A$) de $x\cdot_Ay$, c'est-à-dire $(-x)\cdot_Ay = -(x\cdot_Ay)$. On
montre de façon analogue $x\cdot_A(-y) = -(x\cdot_Ay)$.

(3) En appliquant deux fois le résultat (2) :
\[
(-x)\cdot_A(-y) = -\big(x\cdot_A(-y)\big) = -\big(-(x\cdot_Ay)\big) =
x\cdot_Ay,
\]
la dernière égalité provenant du fait que l'opposé de l'opposé d'un
élément redonne cet élément (propriété générale des groupes).
\end{proof}

\begin{remarque}
Ce théorème redémontre, dans le cadre abstrait des anneaux, la règle
scolaire \og moins par moins donne plus \fg{} : elle n'est pas un axiome
séparé, mais une \textbf{conséquence logique} des axiomes d'anneau
(distributivité, notamment).
\end{remarque}

\subsection{Anneau intègre}


\begin{definition}[4.7]
Un anneau $(A,+_A,\cdot_A)$ est dit \textbf{intègre} si
\[
\forall x,y\in A,\quad \big(x\cdot_A y = 0_A\big) \Rightarrow
\big(x=0_A \ \vee\ y=0_A\big).
\]
Un élément $x\neq0_A$ tel qu'il existe $y\neq0_A$ avec $x\cdot_Ay=0_A$ est
appelé \textbf{diviseur de zéro}.
\end{definition}

\begin{exemple}[4.8]
$(\mathbb{Z},+,\cdot)$, $(\mathbb{Q},+,\cdot)$, $(\mathbb{R},+,\cdot)$,
$(\mathbb{C},+,\cdot)$ sont des anneaux intègres.

En revanche, l'anneau $(\mathbb{Z}/6\mathbb{Z}, \oplus,\odot)$ (avec les
opérations induites, voir Exemple 4.2) \textbf{n'est pas intègre} : on a
$\dot 2 \neq \dot 0$ et $\dot 3 \neq \dot 0$, mais $\dot 2 \odot \dot 3 =
\dot 6 = \dot 0$ (puisque $6$ est congru à $0$ modulo $6$). Les classes
$\dot 2$ et $\dot 3$ sont des diviseurs de zéro dans
$\mathbb{Z}/6\mathbb{Z}$.
\end{exemple}

\section{Structure de corps}


\begin{definition}[4.8]
Soit $(K,+_K,\cdot_K)$ un anneau unitaire. On dit que $(K,+_K,\cdot_K)$ est
un \textbf{corps} si :
\begin{enumerate}
  \item $1_K \neq 0_K$ ;
  \item tout élément de $K\setminus\{0_K\}$ est inversible pour $\cdot_K$.
\end{enumerate}
Le corps est dit \textbf{commutatif} si $\cdot_K$ est commutative.
\end{definition}

\begin{remarque}
Si $(K,+_K,\cdot_K)$ est un corps, alors $(K^*,\cdot_K)$, avec
$K^*=K\setminus\{0_K\}$, est un \textbf{groupe} (c'est exactement la
condition 2 de la définition, combinée à l'associativité et à l'existence
de l'unité héritées de la structure d'anneau).
\end{remarque}

\begin{exemple}[4.9]
$(\mathbb{Q},+,\cdot)$, $(\mathbb{R},+,\cdot)$, $(\mathbb{C},+,\cdot)$ sont
des corps commutatifs. Nous verrons au Chapitre 5 que $(\mathbb{R},+,\cdot)$
et $(\mathbb{C},+,\cdot)$ sont les deux corps centraux de ce cours.

$(\mathbb{Z},+,\cdot)$ \textbf{n'est pas} un corps : les seuls éléments de
$\mathbb{Z}^*$ inversibles pour la multiplication usuelle sont $1$ et
$-1$ (par exemple, $2$ n'a pas d'inverse entier, $\frac12\notin
\mathbb{Z}$).
\end{exemple}

\begin{theoreme}[4.6 -- Tout corps est un anneau intègre]
Si $(K,+_K,\cdot_K)$ est un corps, alors c'est un anneau intègre.
\end{theoreme}

\begin{proof}
Soient $a,b\in K$ tels que $a\cdot_K b = 0_K$. Montrons que $a=0_K$ ou
$b=0_K$, en raisonnant par disjonction de cas sur $a$ (Chapitre 1, section
1.3.4) : ou bien $a=0_K$, et la conclusion est immédiate ; ou bien
$a\neq0_K$, et alors, par définition d'un corps, $a$ est inversible pour
$\cdot_K$. Composons l'égalité $a\cdot_Kb=0_K$ à gauche par $a^{-1}$ :
\[
a^{-1}\cdot_K(a\cdot_Kb) = a^{-1}\cdot_K 0_K
\ \Longrightarrow\
(a^{-1}\cdot_Ka)\cdot_Kb = 0_K \quad\text{(associativité, et Théorème
4.5(1))}
\]
\[
\Longrightarrow\ 1_K\cdot_Kb = 0_K
\ \Longrightarrow\ b=0_K.
\]
Dans les deux cas, $a=0_K$ ou $b=0_K$ : $(K,+_K,\cdot_K)$ est intègre.
\end{proof}

\begin{remarque}[Pourquoi la réciproque est fausse]
Un anneau intègre n'est pas nécessairement un corps : $(\mathbb{Z},+,\cdot)$
est intègre (Exemple 4.8) mais n'est pas un corps (Exemple 4.9), car $2$
n'a pas d'inverse entier. L'intégrité est donc une condition
\textbf{nécessaire mais pas suffisante} pour être un corps.
\end{remarque}


\chapter{Arithmétique dans $\mathbb{Z}$}

\section{Introduction}


Ce chapitre a pour objectifs :
\begin{itemize}
  \item de définir la \textbf{divisibilité} dans $\mathbb{Z}$ et la
  \textbf{division euclidienne} ;
  \item de définir le \textbf{PGCD} de deux entiers et de calculer avec
  l'\textbf{algorithme d'Euclide} ;
  \item de démontrer le \textbf{théorème de Bézout} et le \textbf{théorème
  de Gauss}, et de les appliquer à la résolution d'équations dans
  $\mathbb{Z}$ ;
  \item d'approfondir l'étude des \textbf{congruences} et de la structure
  de $\mathbb{Z}/n\mathbb{Z}$, en caractérisant précisément ses éléments
  inversibles.
\end{itemize}

\begin{application}
L'arithmétique de $\mathbb{Z}$ est le fondement de la \textbf{cryptographie
à clé publique} (RSA repose entièrement sur l'arithmétique modulaire et le
théorème de Bézout), des \textbf{codes correcteurs d'erreurs}, du calcul
des \textbf{clés de contrôle} (IBAN, numéro de sécurité sociale), et de
nombreux \textbf{algorithmes informatiques} (l'algorithme d'Euclide est
l'un des plus anciens algorithmes connus, et reste l'un des plus efficaces
pour calculer un PGCD).
\end{application}

\section{Prérequis}

Ce chapitre s'appuie sur le raisonnement par récurrence et la disjonction
de cas (Chapitre 1), les relations d'équivalence et $\mathbb{Z}/n\mathbb{Z}$
(Chapitre 2), et les structures de groupe, d'anneau et de corps
(Chapitre 4), en particulier la caractérisation des éléments inversibles
d'un anneau.

\section{Divisibilité dans \texorpdfstring{$\mathbb{Z}$}{Z}}

\begin{definition}[5.1]
Soient $a,b\in\mathbb{Z}$. On dit que $a$ \textbf{divise} $b$ (ou que $b$
est un \textbf{multiple} de $a$), noté $a\mid b$, s'il existe
$k\in\mathbb{Z}$ tel que $b=ka$.
\end{definition}

\begin{proposition}[5.1]
Pour tous $a,b,c\in\mathbb{Z}$ :
\begin{enumerate}
  \item $a\mid a$ (réflexivité) ; $1\mid a$ ; $a \mid 0$.
  \item Si $a\mid b$ et $b\mid c$, alors $a\mid c$ (transitivité).
  \item Si $a\mid b$ et $a\mid c$, alors $a \mid (b+c)$ et
  $a \mid (\lambda b + \mu c)$ pour tous $\lambda,\mu\in\mathbb{Z}$.
  \item Si $a\mid b$ et $b\mid a$, alors $a=b$ ou $a=-b$ (i.e. $|a|=|b|$).
\end{enumerate}
\end{proposition}

\begin{proof}[Démonstration de (4)]
Supposons $a\mid b$ et $b\mid a$, avec $a,b$ non nuls (le cas où l'un des
deux est nul est immédiat : si $a\mid b$ et $b=0$, alors $a\mid0$ ce qui
n'apporte rien, mais $b\mid a$ avec $b=0$ force $a=0$, donc $a=b=0$).
Il existe $k,k'\in\mathbb{Z}$ tels que $b=ka$ et $a=k'b$. En substituant :
$a = k'(ka) = (k'k)a$. Comme $a\neq0$, on simplifie : $k'k=1$. Deux
entiers relatifs dont le produit vaut $1$ ne peuvent être que
$(k,k')=(1,1)$ ou $(k,k')=(-1,-1)$ (ce sont les seuls diviseurs entiers de
$1$). Dans le premier cas, $b=a$ ; dans le second, $b=-a$.
\end{proof}

\begin{remarque}
On retrouve ici exactement le même type d'argument que pour l'antisymétrie
de la divisibilité sur $\mathbb{N}$ (Chapitre 2, Théorème 2.2) : sur
$\mathbb{Z}$, la divisibilité n'est \textbf{pas} antisymétrique
(elle donne $|a|=|b|$, pas $a=b$), ce qui confirme qu'elle n'est pas une
relation d'ordre sur $\mathbb{Z}$ tout entier.
\end{remarque}

\section{Division euclidienne}


\begin{theoreme}[5.1 -- Division euclidienne]
Soient $a\in\mathbb{Z}$ et $b\in\mathbb{N}^*$. Il existe un \textbf{unique}
couple $(q,r)\in\mathbb{Z}\times\mathbb{N}$ tel que
\[
a = bq+r \qquad \text{et} \qquad 0\leq r < b.
\]
L'entier $q$ est appelé \textbf{quotient} et $r$ le \textbf{reste} de la
division euclidienne de $a$ par $b$.
\end{theoreme}

\begin{proof}
\textbf{Existence.} Considérons l'ensemble
$S = \{a-bk : k\in\mathbb{Z}\} \cap \mathbb{N}$. Cet ensemble est non
vide : en prenant $k$ suffisamment négatif (par exemple
$k=-|a|$), on a $a-bk = a+b|a| \geq a + |a| \geq 0$ (car $b\geq1$), donc
$S$ contient un élément de $\mathbb{N}$.

Comme $S$ est une partie non vide de $\mathbb{N}$, elle admet un plus
petit élément (propriété fondamentale de $\mathbb{N}$, dite de
\textbf{bon ordre}, admise dans ce cours). Notons $r$ ce plus petit
élément, et $q\in\mathbb{Z}$ tel que $r=a-bq$, c'est-à-dire $a=bq+r$.

Par construction, $r\geq0$. Montrons $r<b$ par l'absurde : si $r\geq b$,
alors $r-b\geq0$, et $r-b = a-bq-b = a-b(q+1) \in S$ (c'est bien de la
forme $a-bk$ avec $k=q+1$). Or $r-b<r$ (car $b>0$), ce qui contredit la
minimalité de $r$ dans $S$. Donc $r<b$.

\textbf{Unicité.} Supposons $a=bq+r=bq'+r'$ avec $0\leq r<b$ et
$0\leq r'<b$. Alors $b(q-q') = r'-r$. Comme $0\leq r<b$ et $0\leq r'<b$,
on a $-b < r'-r < b$, c'est-à-dire $|r'-r|<b$, donc $|b(q-q')|<b$, soit
$|q-q'|<1$. Comme $q-q'\in\mathbb{Z}$, ceci force $q-q'=0$, donc $q=q'$,
puis $r=r'$.
\end{proof}


\begin{exemple}[5.1]
\begin{enumerate}
  \item Division de $47$ par $6$ : $47 = 6\times7+5$ (avec $0\leq5<6$).
  Quotient $q=7$, reste $r=5$.
  \item Division de $-17$ par $5$ : attention, on veut $r\geq0$. On a
  $-17 = 5\times(-4)+3$ (avec $0\leq3<5$). Quotient $q=-4$, reste $r=3$.
  \emph{On ne prend pas $q=-3$ (qui donnerait $-17=5\times(-3)-2$), car
  $-2$ n'est pas dans $[0,5[$.}
\end{enumerate}
\end{exemple}

\section{PGCD et algorithme d'Euclide}


\begin{definition}[5.2]
Soient $a,b\in\mathbb{Z}$, non tous deux nuls. Le \textbf{plus grand
commun diviseur} de $a$ et $b$, noté $\mathrm{pgcd}(a,b)$ (ou $a\wedge b$),
est le plus grand entier naturel qui divise à la fois $a$ et $b$.
\end{definition}

\begin{definition}[5.3]
Deux entiers $a,b$ sont dits \textbf{premiers entre eux} si
$\mathrm{pgcd}(a,b)=1$.
\end{definition}

\begin{proposition}[5.2 -- Principe de l'algorithme d'Euclide]
Soient $a\in\mathbb{Z}$, $b\in\mathbb{N}^*$, et $r$ le reste de la
division euclidienne de $a$ par $b$. Alors
\[
\mathrm{pgcd}(a,b) = \mathrm{pgcd}(b,r).
\]
\end{proposition}

\begin{proof}
Écrivons $a=bq+r$. Montrons que $a$ et $b$ ont exactement les mêmes
diviseurs communs que $b$ et $r$.

Si $d$ divise $a$ et $b$, alors $d$ divise $bq$ (Propriété 5.1(3)) et donc
$d$ divise $a-bq=r$ (Propriété 5.1(3) à nouveau). Donc $d$ divise $b$ et
$r$.

Réciproquement, si $d$ divise $b$ et $r$, alors $d$ divise $bq+r=a$. Donc
$d$ divise $a$ et $b$.

Les diviseurs communs de $(a,b)$ et de $(b,r)$ sont donc exactement les
mêmes ensembles ; en particulier, leurs plus grands éléments coïncident :
$\mathrm{pgcd}(a,b)=\mathrm{pgcd}(b,r)$.
\end{proof}

\begin{theoreme}[5.2 -- Algorithme d'Euclide]
Soient $a,b\in\mathbb{N}^*$. On construit une suite de divisions
euclidiennes successives :
\[
a = bq_1+r_1,\quad b=r_1q_2+r_2,\quad r_1=r_2q_3+r_3,\ \dots
\]
en divisant, à chaque étape, le diviseur précédent par le reste précédent,
jusqu'à obtenir un reste nul. Cette suite de restes $(r_1,r_2,r_3,\dots)$
est \textbf{strictement décroissante} dans $\mathbb{N}$, donc devient nulle
en un nombre fini d'étapes. Le \textbf{dernier reste non nul} est
$\mathrm{pgcd}(a,b)$.
\end{theoreme}

\begin{proof}
La suite $(r_1,r_2,\dots)$ est une suite d'entiers naturels strictement
décroissante ($0\leq r_{i+1}<r_i$ à chaque division euclidienne) : une
telle suite ne peut pas décroître indéfiniment (il n'existe pas de suite
infinie strictement décroissante d'entiers naturels), donc il existe un
rang $n$ tel que $r_n=0$. En appliquant la Propriété 5.2 à chaque étape :
\[
\mathrm{pgcd}(a,b) = \mathrm{pgcd}(b,r_1) = \mathrm{pgcd}(r_1,r_2) = \cdots
= \mathrm{pgcd}(r_{n-2},r_{n-1}) = \mathrm{pgcd}(r_{n-1},0) = r_{n-1},
\]
la dernière égalité provenant du fait que $\mathrm{pgcd}(x,0)=x$ pour tout
$x\in\mathbb{N}^*$ (car $x$ divise $0$, et $x$ est le plus grand diviseur
de $x$ lui-même). Ainsi $\mathrm{pgcd}(a,b)$ est le dernier reste non nul
$r_{n-1}$.
\end{proof}

\begin{exemple}[5.2]
Calculons $\mathrm{pgcd}(126,35)$ par l'algorithme d'Euclide :
\[
126 = 3\times35+21, \qquad 35=1\times21+14, \qquad 21=1\times14+7, \qquad
14=2\times7+0.
\]
Le dernier reste non nul est $7$ : $\mathrm{pgcd}(126,35)=7$.
\end{exemple}

\section{Théorème de Bézout}


\begin{theoreme}[5.3 -- Théorème de Bézout]
Soient $a,b\in\mathbb{Z}$, non tous deux nuls, et $d=\mathrm{pgcd}(a,b)$.
Il existe $u,v\in\mathbb{Z}$ tels que
\[
au+bv = d.
\]
\end{theoreme}

\begin{proof}
Considérons l'ensemble $S=\{au+bv : u,v\in\mathbb{Z}\} \cap \mathbb{N}^*$
(les combinaisons entières de $a$ et $b$ qui sont strictement positives).
$S$ est non vide : par exemple, $a^2+b^2 = a\times a+b\times b \in S$ (à
moins que $a=b=0$, exclu par hypothèse), et $a^2+b^2>0$.

Comme $S\subset\mathbb{N}^*$ est non vide, il admet un plus petit élément,
noté $\delta$, avec $\delta=au_0+bv_0$ pour certains $u_0,v_0\in\mathbb{Z}$.

\textbf{Montrons que $\delta$ divise $a$.} Effectuons la division
euclidienne de $a$ par $\delta$ : $a=\delta q+r$ avec $0\leq r<\delta$.
Alors
\[
r = a-\delta q = a - (au_0+bv_0)q = a(1-u_0q) + b(-v_0q),
\]
qui est bien une combinaison entière de $a$ et $b$. Si $r>0$, alors
$r\in S$ et $r<\delta$, ce qui contredit la minimalité de $\delta$ dans
$S$. Donc $r=0$, c'est-à-dire $\delta$ divise $a$. On montre de façon
identique que $\delta$ divise $b$.

$\delta$ est donc un diviseur commun de $a$ et $b$ : par définition du
PGCD, $\delta \leq d$.

\textbf{Montrons que $d$ divise $\delta$.} Comme $d$ divise $a$ et $b$, il
divise toute combinaison entière $au+bv$ (Propriété 5.1(3)), en
particulier $\delta=au_0+bv_0$. Donc $d$ divise $\delta$, et comme
$\delta>0$, ceci entraîne $d\leq\delta$.

En combinant $\delta\leq d$ et $d\leq\delta$, on obtient $\delta=d$. Ainsi
$d = au_0+bv_0$ : le couple $(u,v)=(u_0,v_0)$ convient.
\end{proof}

\begin{remarque}[Comment trouver $u,v$ en pratique]
La démonstration ci-dessus prouve l'\emph{existence} de $u,v$, mais ne
donne pas de méthode de calcul directe. En pratique, on retrouve $u,v$
en \textbf{remontant} les étapes de l'algorithme d'Euclide, en exprimant
à chaque étape le reste en fonction des deux précédents.
\end{remarque}

\begin{exemple}[5.3]
Trouvons $u,v\in\mathbb{Z}$ tels que $126u+35v=7=\mathrm{pgcd}(126,35)$
(voir Exemple 5.2 pour les divisions). On \textbf{remonte} les calculs, en
partant de l'avant-dernière ligne :
\[
7 = 21 - 1\times14 \qquad (\text{issu de } 21=1\times14+7)
\]
On remplace $14$ par son expression $14=35-1\times21$ :
\[
7 = 21 - 1\times(35-1\times21) = 2\times21 - 35.
\]
On remplace $21$ par son expression $21=126-3\times35$ :
\[
7 = 2\times(126-3\times35) - 35 = 2\times126 - 6\times35 - 35 =
2\times126 - 7\times35.
\]
On obtient $126\times2 + 35\times(-7) = 252-245=7$ : on vérifie que
$u=2$, $v=-7$ conviennent.
\end{exemple}

\begin{theoreme}[5.4 -- Corollaire : caractérisation des nombres premiers entre eux]
Deux entiers $a,b\in\mathbb{Z}$ sont premiers entre eux si et seulement si
il existe $u,v\in\mathbb{Z}$ tels que $au+bv=1$.
\end{theoreme}

\begin{proof}
\textbf{Sens direct.} Si $\mathrm{pgcd}(a,b)=1$, le Théorème 5.3 (appliqué
avec $d=1$) donne directement l'existence de $u,v$ tels que $au+bv=1$.

\textbf{Sens réciproque.} Supposons qu'il existe $u,v$ avec $au+bv=1$.
Soit $d=\mathrm{pgcd}(a,b)$. Comme $d$ divise $a$ et $b$, il divise
$au+bv=1$ (Propriété 5.1(3)). Un entier naturel qui divise $1$ ne peut
être que $1$ (car $d\geq1$ et $d\mid1$ implique $d\leq1$). Donc $d=1$ :
$a$ et $b$ sont premiers entre eux.
\end{proof}

\begin{theoreme}[5.5 -- Théorème de Gauss]
Soient $a,b,c\in\mathbb{Z}$. Si $a\mid bc$ et $\mathrm{pgcd}(a,b)=1$, alors
$a\mid c$.
\end{theoreme}

\begin{proof}
Comme $\mathrm{pgcd}(a,b)=1$, le Théorème 5.4 donne $u,v\in\mathbb{Z}$
tels que $au+bv=1$. Multiplions cette égalité par $c$ :
\[
auc + bvc = c.
\]
Le premier terme $auc = a(uc)$ est divisible par $a$. Le second terme
$bvc = (bc)v$ est divisible par $a$, car $a\mid bc$ par hypothèse (donc
$a$ divise tout multiple de $bc$, en particulier $(bc)v$). Une somme de
deux multiples de $a$ est un multiple de $a$ (Propriété 5.1(3)) : donc $a$
divise $auc+bvc=c$.
\end{proof}


\section{Congruences et structure de \texorpdfstring{$\mathbb{Z}/n\mathbb{Z}$}{Z/nZ}}


\begin{remarque}[Rappel]
La relation de congruence modulo $n$ ($n\in\mathbb{N}^*$) sur
$\mathbb{Z}$ est définie par $x\equiv y \pmod n \iff n \mid (x-y)$
(Chapitre 2, Exemple 2.8) ; c'est une relation d'équivalence, et l'ensemble
quotient $\mathbb{Z}/n\mathbb{Z}$, muni des opérations $\oplus$ et
$\odot$ induites, est un anneau commutatif unitaire (Chapitre 4).
\end{remarque}

\begin{theoreme}[5.6 -- Éléments inversibles de \texorpdfstring{$\mathbb{Z}/n\mathbb{Z}$}{Z/nZ}]
Soit $n\in\mathbb{N}$, $n\geq2$, et $a\in\mathbb{Z}$. La classe
$\dot a \in \mathbb{Z}/n\mathbb{Z}$ est inversible pour $\odot$ si et
seulement si $a$ et $n$ sont premiers entre eux.
\end{theoreme}

\begin{proof}
\textbf{Sens direct.} Supposons $\dot a$ inversible : il existe
$u\in\mathbb{Z}$ tel que $\dot a \odot \dot u = \dot 1$, c'est-à-dire
$au \equiv 1 \pmod n$, c'est-à-dire $n \mid (au-1)$. Il existe donc
$k\in\mathbb{Z}$ tel que $au-1=nk$, soit $au-nk=1$ (avec
$v=-k\in\mathbb{Z}$ : $au+nv=1$). D'après le Théorème 5.4,
$\mathrm{pgcd}(a,n)=1$.

\textbf{Sens réciproque.} Supposons $\mathrm{pgcd}(a,n)=1$. D'après le
Théorème 5.3 (Bézout), il existe $u,v\in\mathbb{Z}$ tels que $au+nv=1$.
En passant aux classes modulo $n$ : $\dot a \odot \dot u \oplus \dot n
\odot \dot v = \dot 1$. Or $\dot n = \dot 0$ (car $n$ est congru à $0$
modulo $n$), donc $\dot n \odot \dot v = \dot 0$. On obtient
$\dot a \odot \dot u = \dot 1$ : $\dot a$ est inversible, d'inverse
$\dot u$.
\end{proof}

\begin{exemple}[5.4]
Déterminons l'inverse de $17$ dans $\mathbb{Z}/40\mathbb{Z}$. Par
l'algorithme d'Euclide étendu (calcul analogue à l'Exemple 5.3), on
trouve $3\times40 - 7\times17 = 1$, c'est-à-dire $17\times(-7) + 40\times3
= 1$. Donc $\dot{17} \odot \dot{(-7)} = \dot 1$ dans
$\mathbb{Z}/40\mathbb{Z}$ : l'inverse de $\dot{17}$ est $\dot{(-7)} =
\dot{33}$ (car $-7+40=33$).

Vérification : $17\times33 = 561 = 14\times40+1$, donc $561 \equiv 1
\pmod{40}$. Correct.
\end{exemple}

\begin{theoreme}[5.7 -- Corollaire : caractérisation des corps \texorpdfstring{$\mathbb{Z}/n\mathbb{Z}$}{Z/nZ}]
Soit $n\in\mathbb{N}$, $n\geq2$. L'anneau $(\mathbb{Z}/n\mathbb{Z},
\oplus,\odot)$ est un \textbf{corps} si et seulement si $n$ est
\textbf{premier}.
\end{theoreme}

\begin{proof}
\textbf{Si $n$ est premier.} Pour tout $a\in\{1,\dots,n-1\}$, comme $n$
est premier, les seuls diviseurs de $n$ sont $1$ et $n$ ; comme
$1\leq a<n$, $a$ n'est pas un multiple de $n$, donc le seul diviseur
commun possible à $a$ et $n$ est $1$ : $\mathrm{pgcd}(a,n)=1$. D'après le
Théorème 5.6, $\dot a$ est donc inversible, pour tout
$a\in\{1,\dots,n-1\}$, c'est-à-dire pour tout élément non nul de
$\mathbb{Z}/n\mathbb{Z}$. Comme $\dot0\neq\dot1$ (car $n\geq2$),
$(\mathbb{Z}/n\mathbb{Z},\oplus,\odot)$ est bien un corps.

\textbf{Si $n$ n'est pas premier.} C'est exactement le contenu de
l'Exercice D5 du Chapitre 4 : on écrit $n=ab$ avec $1<a,b<n$, et on montre
que $\dot a \odot \dot b = \dot 0$ avec $\dot a,\dot b\neq\dot 0$ :
l'anneau n'est pas intègre, donc (Théorème 4.6, par contraposée) ce n'est
pas un corps.
\end{proof}

\begin{remarque}
Ce théorème boucle élégamment la boucle ouverte au Chapitre 4 : nous avions
vérifié \og à la main \fg{} que $\mathbb{Z}/5\mathbb{Z}\setminus\{\dot0\}$
est un groupe multiplicatif (car $5$ est premier), et que
$\mathbb{Z}/6\mathbb{Z}$ n'est pas intègre (car $6=2\times3$ n'est pas
premier). Le Théorème 5.7 généralise ces observations ponctuelles en un
critère complet et définitif.
\end{remarque}


\chapter{Nombres complexes}

\section{Introduction}


Ce chapitre a pour objectifs :
\begin{itemize}
  \item de \textbf{construire} rigoureusement $\mathbb{C}$ comme corps
  commutatif, et de manipuler sa \textbf{forme algébrique} $a+ib$ ;
  \item d'interpréter géométriquement les nombres complexes dans le
  \textbf{plan complexe} (affixe, module, argument) ;
  \item d'utiliser les \textbf{formes trigonométrique et exponentielle},
  la \textbf{formule d'Euler} et la \textbf{formule de Moivre} ;
  \item de calculer les \textbf{racines $n$-ièmes} d'un nombre complexe ;
  \item de pratiquer la \textbf{linéarisation} de $\cos^n\theta$,
  $\sin^n\theta$, technique essentielle pour le calcul intégral en
  analyse.
\end{itemize}

\begin{application}
Les nombres complexes sont essentiels en \textbf{électronique} (analyse
des circuits en courant alternatif, impédances complexes), en
\textbf{traitement du signal} (transformée de Fourier), en
\textbf{mécanique quantique} (fonction d'onde), et en \textbf{infographie}
(rotations et transformations du plan représentées par la multiplication
complexe).
\end{application}

\section{Prérequis}

Ce chapitre s'appuie sur les structures de corps (Chapitre 4), sur les
identités trigonométriques usuelles ($\cos(a+b)$, $\sin(a+b)$, etc.,
supposées connues du secondaire), et sur la représentation du plan muni
d'un repère orthonormé.

\section{Construction du corps des nombres complexes}

\begin{definition}
On définit sur $\mathbb{R}^2$ l'addition $+$ et la multiplication $\cdot$
par :
\[
(a_1,b_1) + (a_2,b_2) = (a_1+a_2,\ b_1+b_2), \qquad
(a_1,b_1) \cdot (a_2,b_2) = (a_1a_2-b_1b_2,\ a_1b_2+a_2b_1).
\]
\end{definition}

\begin{theoreme}
$(\mathbb{R}^2,+,\cdot)$ est un \textbf{corps commutatif}, appelé
\textbf{corps des nombres complexes} et noté $(\mathbb{C},+,\cdot)$.
\end{theoreme}

\begin{proof}[Démonstration (points essentiels)]
La plupart des axiomes d'anneau commutatif (associativité, commutativité,
distributivité, éléments neutres $(0,0)$ et $(1,0)$) se vérifient par un
calcul direct et immédiat sur les formules de la Définition 6.1, de façon
analogue à ce qui a été fait pour $\mathbb{R}$ ; nous ne le détaillons pas
ici. Le point \textbf{non trivial} est l'existence d'un \textbf{inverse}
pour tout élément non nul, que nous démontrons complètement.

Soit $z=(a,b) \neq (0,0)$ (donc $a^2+b^2\neq0$). Cherchons
$z'=(a',b')$ tel que $z\cdot z' = (1,0)$ :
\[
(aa'-bb',\ ab'+a'b) = (1,0) \iff
\begin{cases} aa'-bb' = 1 \\ ba'+ab' = 0. \end{cases}
\]
C'est un système linéaire de deux équations à deux inconnues $(a',b')$,
de déterminant $a\times a - (-b)\times b = a^2+b^2 \neq 0$ (car
$z\neq(0,0)$). Le système admet donc une unique solution, que l'on
calcule directement :
\[
a' = \frac{a}{a^2+b^2}, \qquad b' = \frac{-b}{a^2+b^2}.
\]
On vérifie que ce couple convient bien (calcul direct de substitution).
Donc tout $z\neq(0,0)$ est inversible : $(\mathbb{C}^*,\cdot)$ est un
groupe, et $1_\mathbb{C}=(1,0)\neq(0,0)_\mathbb{C}$ : $(\mathbb{C},+,\cdot)$
est bien un corps.
\end{proof}

\begin{definition}[6.2 -- Forme algébrique]
En identifiant $(a,0)$ à $a$ (ce qui plonge $\mathbb{R}$ dans
$\mathbb{C}$, $\mathbb{R}\subset\mathbb{C}$) et en notant $i=(0,1)$, on a :
\[
i^2 = (0,1)\cdot(0,1) = (0\times0-1\times1,\ 0\times1+1\times0) = (-1,0)
= -1.
\]
Tout élément $z=(a,b)\in\mathbb{C}$ s'écrit alors de façon
\textbf{unique} :
\[
z = a+ib, \qquad a,b\in\mathbb{R}.
\]
Cette écriture est la \textbf{forme algébrique} de $z$. Le réel $a$ est la
\textbf{partie réelle} de $z$, notée $\mathrm{Re}(z)$ ; le réel $b$ est la
\textbf{partie imaginaire} de $z$, notée $\mathrm{Im}(z)$ (c'est un
\textbf{réel}, pas un complexe !). Si $\mathrm{Im}(z)=0$, $z$ est
\textbf{réel} ; si $\mathrm{Re}(z)=0$ et $z\neq0$, $z$ est dit
\textbf{imaginaire pur}.
\end{definition}


\begin{definition}[Conjugué et module]
Soit $z=a+ib\in\mathbb{C}$ ($a,b\in\mathbb{R}$).
\begin{itemize}
  \item Le \textbf{conjugué} de $z$ est $\bar z = a-ib$.
  \item Le \textbf{module} de $z$ est $|z| = \sqrt{a^2+b^2} \in
  \mathbb{R}^+$.
\end{itemize}
\end{definition}

\begin{proposition}[6.1]
Pour tout $z\in\mathbb{C}$ :
\[
z\bar z = |z|^2, \qquad \mathrm{Re}(z) = \frac{z+\bar z}{2}, \qquad
\mathrm{Im}(z) = \frac{z-\bar z}{2i}.
\]
Si $z\neq0$, l'inverse de $z$ est $z^{-1} = \dfrac{\bar z}{|z|^2}$.
\end{proposition}

\begin{proof}
Écrivons $z=a+ib$. Alors
\[
z\bar z = (a+ib)(a-ib) = a^2 - (ib)^2 = a^2 - i^2b^2 = a^2+b^2 = |z|^2
\]
(en utilisant $i^2=-1$). Les deux autres formules s'obtiennent en
résolvant le système $z=a+ib,\ \bar z=a-ib$ (par somme et différence).
Enfin, si $z\neq0$, $|z|^2\neq0$, et $z\times\dfrac{\bar z}{|z|^2} =
\dfrac{z\bar z}{|z|^2} = \dfrac{|z|^2}{|z|^2}=1$ : on retrouve la formule
d'inversion établie au Théorème 6.1.
\end{proof}

\begin{exemple}[6.1]
Mettons sous forme algébrique $z=(3+2i)(1-i)$ :
\[
z = 3\times1 + 3\times(-i) + 2i\times1 + 2i\times(-i) = 3-3i+2i-2i^2 =
3-i+2 = 5-i
\]
(en utilisant $i^2=-1$, donc $-2i^2=2$).
\end{exemple}

\section{Représentation géométrique}


\begin{definition}[6.4]
Le plan muni d'un repère orthonormé $(O,\vec u,\vec v)$ est appelé
\textbf{plan complexe}. Au nombre complexe $z=a+ib$, on associe le point
$M(a,b)$, appelé \textbf{image} de $z$, et $z$ est appelé
l'\textbf{affixe} de $M$. Le module $|z|$ est alors la \textbf{distance}
$OM$, et l'\textbf{argument} de $z$ (pour $z\neq0$), noté $\arg(z)$, est
une mesure de l'angle $(\vec u, \overrightarrow{OM})$, définie modulo
$2\pi$.
\end{definition}

\begin{center}
\begin{figure}[h]
  \centering
  \includegraphics[width=0.6\textwidth]{image3}
  %\caption{Graphe de la fonction $f$}
\end{figure}
\end{center}

\section{Formes trigonométrique et exponentielle}


\begin{definition}[6.5]
Pour $z\in\mathbb{C}^*$, de module $r=|z|$ et d'argument $\theta=\arg(z)$,
on a $z = r(\cos\theta+i\sin\theta)$ : c'est la \textbf{forme
trigonométrique} de $z$.
\end{definition}

\begin{definition}[6.6 -- Formule d'Euler, forme exponentielle]
Pour $\theta\in\mathbb{R}$, on pose, par définition :
\[
e^{i\theta} := \cos\theta + i\sin\theta \qquad (\textbf{formule d'Euler}).
\]
Tout $z\in\mathbb{C}^*$ s'écrit alors $z=re^{i\theta}$, avec $r=|z|>0$ et
$\theta=\arg(z)$ : c'est la \textbf{forme exponentielle} de $z$.
\end{definition}

\begin{theoreme}[6.2 -- Règles de calcul sur la forme exponentielle]
Pour tous $\theta,\theta'\in\mathbb{R}$ et $n\in\mathbb{Z}$ :
\[
e^{i\theta}\times e^{i\theta'} = e^{i(\theta+\theta')}, \qquad
\overline{e^{i\theta}} = e^{-i\theta} = \frac{1}{e^{i\theta}}, \qquad
\big(e^{i\theta}\big)^n = e^{in\theta} \quad \textbf{(formule de Moivre)}.
\]
\end{theoreme}

\begin{proof}
\textbf{Produit.} En utilisant la Définition 6.6 et les formules
trigonométriques d'addition (rappelées en prérequis) :
\[
e^{i\theta}e^{i\theta'} = (\cos\theta+i\sin\theta)(\cos\theta'+i\sin\theta')
\]
\[
= (\cos\theta\cos\theta'-\sin\theta\sin\theta') +
i(\cos\theta\sin\theta'+\sin\theta\cos\theta')
= \cos(\theta+\theta') + i\sin(\theta+\theta') = e^{i(\theta+\theta')},
\]
en reconnaissant les formules d'addition $\cos(\theta+\theta')$ et
$\sin(\theta+\theta')$.

\textbf{Conjugué et inverse.} $\overline{e^{i\theta}} =
\overline{\cos\theta+i\sin\theta} = \cos\theta-i\sin\theta =
\cos(-\theta)+i\sin(-\theta) = e^{-i\theta}$. Comme
$|e^{i\theta}|=\sqrt{\cos^2\theta+\sin^2\theta}=1$, on a
$e^{i\theta}\times\overline{e^{i\theta}} = |e^{i\theta}|^2 = 1$
(Propriété 6.1), donc $\dfrac{1}{e^{i\theta}} = \overline{e^{i\theta}} =
e^{-i\theta}$.

\textbf{Formule de Moivre.} Par récurrence sur $n\in\mathbb{N}$, en
utilisant la règle du produit établie ci-dessus :
\emph{Initialisation} ($n=0$) : $(e^{i\theta})^0=1=e^{i\times0}$.
\emph{Hérédité} : si $(e^{i\theta})^n=e^{in\theta}$, alors
$(e^{i\theta})^{n+1} = (e^{i\theta})^n \times e^{i\theta} = e^{in\theta}
\times e^{i\theta} = e^{i(n+1)\theta}$. La formule s'étend aux entiers
négatifs grâce à la règle de l'inverse.
\end{proof}

\begin{remarque}[Formule de Moivre -- forme trigonométrique]
En écrivant $e^{i\theta}=\cos\theta+i\sin\theta$, la formule de Moivre
s'énonce aussi :
\[
(\cos\theta+i\sin\theta)^n = \cos(n\theta)+i\sin(n\theta), \qquad
n\in\mathbb{Z}.
\]
\end{remarque}

\begin{exemple}[6.2]
Écrivons $z=1+i$ sous forme trigonométrique. Le module est $|z| =
\sqrt{1^2+1^2}=\sqrt2$. Pour l'argument : $\cos\theta =
\dfrac{1}{\sqrt2} = \dfrac{\sqrt2}{2}$ et $\sin\theta=\dfrac{1}{\sqrt2}=
\dfrac{\sqrt2}{2}$, d'où $\theta=\dfrac{\pi}{4}$. Donc :
\[
z = \sqrt2\left(\cos\frac{\pi}{4}+i\sin\frac{\pi}{4}\right) = \sqrt2\,
e^{i\pi/4}.
\]
\end{exemple}

\begin{exemple}[6.3]
Calculons $(1+i)^{10}$ à l'aide de la forme exponentielle (Exemple 6.2) et
de la formule de Moivre :
\[
(1+i)^{10} = \left(\sqrt2\,e^{i\pi/4}\right)^{10} = (\sqrt2)^{10}\,
e^{i10\pi/4} = 2^5\, e^{i5\pi/2}.
\]
Comme $\dfrac{5\pi}{2} = 2\pi+\dfrac{\pi}{2}$, on a $e^{i5\pi/2} =
e^{i\pi/2} = i$ (l'argument est défini modulo $2\pi$). Donc
$(1+i)^{10} = 32\times i = 32i$.
\end{exemple}

\section{Racines \texorpdfstring{$n$}{n}-ièmes d'un nombre complexe}


\begin{theoreme}[6.3 -- Racines \texorpdfstring{$n$}{n}-ièmes]
Soit $z_0 = r_0 e^{i\theta_0} \in \mathbb{C}^*$ et $n\in\mathbb{N}^*$.
L'équation $z^n=z_0$ admet exactement $n$ solutions distinctes dans
$\mathbb{C}$, données par :
\[
z_k = r_0^{1/n}\, e^{i\left(\frac{\theta_0}{n}+\frac{2k\pi}{n}\right)},
\qquad k=0,1,\dots,n-1.
\]
\end{theoreme}

\begin{proof}
Cherchons $z=re^{i\theta}$ (avec $r>0$) tel que $z^n=z_0$. D'après la
formule de Moivre, $z^n = r^n e^{in\theta}$. L'égalité $z^n=z_0$
équivaut donc à
\[
r^n e^{in\theta} = r_0 e^{i\theta_0}.
\]
En égalant les modules (qui sont des réels positifs, donc l'égalité de
deux nombres complexes écrits sous forme exponentielle avec modules
positifs équivaut à l'égalité des modules et des arguments modulo
$2\pi$) : $r^n = r_0$, donc $r = r_0^{1/n}$ (racine $n$-ième réelle
positive, qui existe et est unique car $r_0>0$).

Pour les arguments : $n\theta \equiv \theta_0 \pmod{2\pi}$, c'est-à-dire
$n\theta = \theta_0+2k\pi$ pour un certain $k\in\mathbb{Z}$, soit
$\theta = \dfrac{\theta_0}{n}+\dfrac{2k\pi}{n}$.

Il reste à vérifier que les valeurs de $k$ donnant des solutions
\textbf{distinctes} sont exactement $k=0,1,\dots,n-1$ : deux valeurs
$k$ et $k'$ donnent le même $\theta$ modulo $2\pi$ si et seulement si
$\dfrac{2(k-k')\pi}{n} \equiv 0 \pmod{2\pi}$, c'est-à-dire si et
seulement si $n$ divise $k-k'$. Les entiers $0,1,\dots,n-1$ sont
précisément un représentant de chaque classe modulo $n$ (Chapitre 2) :
ils donnent donc $n$ arguments deux à deux distincts modulo $2\pi$,
c'est-à-dire $n$ solutions $z_k$ deux à deux distinctes.
\end{proof}

\begin{definition}[6.7 -- Racines \texorpdfstring{$n$}{n}-ièmes de l'unité]
Les solutions de $z^n=1$ sont appelées \textbf{racines $n$-ièmes de
l'unité}. D'après le Théorème 6.3 (avec $r_0=1,\theta_0=0$), ce sont les
$n$ nombres complexes
\[
\omega_k = e^{i2k\pi/n}, \qquad k=0,1,\dots,n-1.
\]
\end{definition}

\begin{exemple}[6.4]
Résolvons $z^3=1$ (racines cubiques de l'unité) : $\omega_k =
e^{i2k\pi/3}$ pour $k=0,1,2$, soit :
\[
\omega_0 = 1, \qquad \omega_1 = e^{i2\pi/3} = -\frac12+i\frac{\sqrt3}{2},
\qquad \omega_2 = e^{i4\pi/3} = -\frac12-i\frac{\sqrt3}{2}.
\]
Géométriquement, ces trois points forment un triangle équilatéral inscrit
dans le cercle unité.
\end{exemple}

\begin{proposition}[6.2 -- Somme des racines \texorpdfstring{$n$}{n}-ièmes de l'unité]
Pour $n\geq2$, la somme des racines $n$-ièmes de l'unité est nulle :
\[
\sum_{k=0}^{n-1} \omega_k = 0.
\]
\end{proposition}

\begin{proof}
On reconnaît une somme géométrique de raison $\omega_1 = e^{i2\pi/n}
\neq 1$ (car $n\geq2$) :
\[
\sum_{k=0}^{n-1} \omega_k = \sum_{k=0}^{n-1} \big(\omega_1\big)^k =
\frac{\omega_1^n - 1}{\omega_1 - 1}
\]
(formule de la somme géométrique, valable pour $\omega_1\neq1$). Or
$\omega_1^n = e^{i2\pi} = 1$ : le numérateur est nul. Comme le
dénominateur $\omega_1-1\neq0$, toute la fraction est nulle.
\end{proof}

\section{Linéarisation et formule d'Euler}


\begin{proposition}[6.3 -- Formules d'Euler]
Pour $\theta\in\mathbb{R}$ :
\[
\cos\theta = \frac{e^{i\theta}+e^{-i\theta}}{2}, \qquad
\sin\theta = \frac{e^{i\theta}-e^{-i\theta}}{2i}.
\]
\end{proposition}

\begin{proof}
D'après la Définition 6.6, $e^{i\theta}=\cos\theta+i\sin\theta$ et
$e^{-i\theta}=\cos\theta-i\sin\theta$. En additionnant ces deux
égalités : $e^{i\theta}+e^{-i\theta} = 2\cos\theta$, d'où la première
formule. En les soustrayant : $e^{i\theta}-e^{-i\theta} = 2i\sin\theta$,
d'où la seconde.
\end{proof}

\begin{exemple}[6.5 -- Linéarisation de \texorpdfstring{$\cos^2\theta$}{cos²θ}]
\[
\cos^2\theta = \left(\frac{e^{i\theta}+e^{-i\theta}}{2}\right)^2 =
\frac{e^{2i\theta}+2+e^{-2i\theta}}{4} =
\frac{e^{2i\theta}+e^{-2i\theta}}{4} + \frac12
= \frac{2\cos(2\theta)}{4}+\frac12 = \frac{1+\cos(2\theta)}{2},
\]
en reconnaissant $e^{2i\theta}+e^{-2i\theta}=2\cos(2\theta)$ (Propriété
6.3 appliquée à l'angle $2\theta$).
\end{exemple}

\begin{exemple}[6.6 -- Linéarisation de \texorpdfstring{$\sin^3\theta$}{sin³θ}]
\[
\sin^3\theta = \left(\frac{e^{i\theta}-e^{-i\theta}}{2i}\right)^3 =
\frac{(e^{i\theta}-e^{-i\theta})^3}{-8i}.
\]
Développons le numérateur par la formule du binôme :
\[
(e^{i\theta}-e^{-i\theta})^3 = e^{3i\theta} - 3e^{2i\theta}e^{-i\theta} +
3e^{i\theta}e^{-2i\theta} - e^{-3i\theta}
= e^{3i\theta} - 3e^{i\theta} + 3e^{-i\theta} - e^{-3i\theta}
\]
\[
= \big(e^{3i\theta}-e^{-3i\theta}\big) - 3\big(e^{i\theta}-e^{-i\theta}\big)
= 2i\sin(3\theta) - 3\times2i\sin\theta = 2i\big(\sin(3\theta)-3\sin\theta\big).
\]
En reportant :
\[
\sin^3\theta = \frac{2i\big(\sin(3\theta)-3\sin\theta\big)}{-8i} =
\frac{3\sin\theta - \sin(3\theta)}{4}.
\]
\end{exemple}

\begin{remarque}[Opération inverse : développer \texorpdfstring{$\cos(n\theta)$}{cos(nθ)}]
La même formule de Moivre permet aussi l'opération \textbf{inverse} :
développer $\cos(n\theta)$ en fonction de $\cos\theta$. Par exemple, la
formule de Moivre à l'ordre $3$ donne $(\cos\theta+i\sin\theta)^3 =
\cos(3\theta)+i\sin(3\theta)$. En développant le membre de gauche par le
binôme et en identifiant les parties réelles :
\[
\cos^3\theta - 3\cos\theta\sin^2\theta = \cos(3\theta),
\]
puis, en remplaçant $\sin^2\theta=1-\cos^2\theta$ :
\[
\cos(3\theta) = \cos^3\theta - 3\cos\theta(1-\cos^2\theta) = 4\cos^3\theta
- 3\cos\theta.
\]
\end{remarque}


\chapter{Polynômes et fractions rationnelles}

\section{Introduction}


Ce chapitre a pour objectifs :
\begin{itemize}
  \item de définir rigoureusement un \textbf{polynôme}, son \textbf{degré}
  et les opérations sur $K[X]$ ;
  \item d'étudier les \textbf{racines} d'un polynôme et leur \textbf{ordre
  de multiplicité} ;
  \item de maîtriser la \textbf{division euclidienne} et la \textbf{division
  suivant les puissances croissantes} dans $K[X]$ ;
  \item de connaître la \textbf{factorisation} des polynômes dans
  $\mathbb{C}[X]$ et $\mathbb{R}[X]$ ;
  \item de \textbf{décomposer une fraction rationnelle en éléments
  simples}.
\end{itemize}

\begin{application}
Les polynômes interviennent dans l'\textbf{interpolation numérique}
(construire une courbe passant par des points donnés), les
\textbf{codes correcteurs d'erreurs} (codes de Reed-Solomon, utilisés dans
les CD, DVD, QR codes), le \textbf{calcul formel}, et la décomposition en
éléments simples est l'outil de base du calcul de \textbf{primitives de
fractions rationnelles} en analyse, ainsi que de la \textbf{transformée en
$Z$} en traitement du signal.
\end{application}

\section{Prérequis}

Ce chapitre s'appuie sur les structures d'anneau et de corps (Chapitre 4),
sur l'arithmétique de $\mathbb{Z}$ -- l'analogie avec $K[X]$ est
omniprésente (Chapitre 5) --, et sur les nombres complexes (Chapitre 6),
en particulier pour la factorisation des polynômes.

\section{Polynômes à une indéterminée}

\begin{definition}[7.1]
Soit $K$ un corps commutatif ($K=\mathbb{R}$ ou $K=\mathbb{C}$ dans ce
cours). On appelle \textbf{polynôme} à une indéterminée $X$, à
coefficients dans $K$, toute écriture
\[
P = a_0 + a_1X + a_2X^2 + \cdots + a_nX^n + \cdots
\]
où les $a_i \in K$ (appelés \textbf{coefficients} de $P$) sont nuls sauf
un nombre fini d'entre eux. L'ensemble des polynômes à coefficients dans
$K$ est noté $K[X]$.
\begin{itemize}
  \item Le plus grand indice $n$ tel que $a_n\neq0$ (s'il existe) est
  appelé \textbf{degré} de $P$, noté $\deg P$, et $a_nX^n$ est le
  \textbf{terme dominant} de $P$.
  \item Si tous les $a_i$ sont nuls, $P$ est le \textbf{polynôme nul},
  noté $0$ ; par convention $\deg 0 = -\infty$.
\end{itemize}
\end{definition}

\begin{definition}[7.2 -- Opérations]
Soient $P=\sum a_iX^i$ et $Q=\sum b_iX^i$ deux polynômes de $K[X]$. On
définit :
\[
P+Q = \sum_i (a_i+b_i)X^i, \qquad P\cdot Q = \sum_n \left(\sum_{i+j=n}
a_ib_j\right) X^n.
\]
\end{definition}

\begin{theoreme}[7.1]
$(K[X],+,\cdot)$ est un \textbf{anneau commutatif unitaire}. De plus, pour
tous $P,Q\in K[X]$ :
\[
\deg(P+Q) \leq \max(\deg P,\deg Q), \qquad \deg(P\cdot Q) = \deg P + \deg Q.
\]
En particulier, $K[X]$ est \textbf{intègre}.
\end{theoreme}

\begin{proof}[Démonstration (formule du degré du produit)]
La structure d'anneau se vérifie par un calcul direct sur les formules de
la Définition 7.2, de façon analogue à ce qui a été fait dans
$\mathbb{Z}$ ou $\mathbb{C}$ ; nous admettons ce point et détaillons
plutôt la formule du degré, plus spécifique.

Posons $n=\deg P$, $m=\deg Q$ (supposons $P,Q\neq0$, sinon le résultat est
immédiat par convention sur $-\infty$). Le coefficient de $X^{n+m}$ dans
$P\cdot Q$ est
\[
\sum_{i+j=n+m} a_ib_j = a_nb_m + \underbrace{\sum_{\substack{i+j=n+m \\
i<n}} a_ib_j}_{=0 \text{ car alors } j>m,\ b_j=0} +
\underbrace{\sum_{\substack{i+j=n+m \\ i>n}} a_ib_j}_{=0 \text{ car alors
} a_i=0}
= a_nb_m.
\]
Comme $K$ est un \textbf{corps} (donc intègre, Théorème 4.6), $a_n\neq0$
et $b_m\neq0$ entraînent $a_nb_m\neq0$. Donc le coefficient de $X^{n+m}$
dans $P\cdot Q$ est non nul, et on vérifie (par un argument analogue) que
tous les coefficients de $X^k$ pour $k>n+m$ sont nuls. Donc
$\deg(PQ)=n+m=\deg P+\deg Q$.

\textbf{Intégrité :} si $P\cdot Q=0$, alors $\deg(PQ)=-\infty$, donc
$\deg P+\deg Q=-\infty$, ce qui force $\deg P=-\infty$ ou $\deg
Q=-\infty$ (par convention sur les degrés), c'est-à-dire $P=0$ ou $Q=0$.
\end{proof}


\section{Racines d'un polynôme}

\begin{definition}[7.3]
Soit $P=a_0+a_1X+\cdots+a_nX^n \in K[X]$. On appelle \textbf{fonction
polynôme} associée à $P$ la fonction $\widetilde P : K\to K$,
$\widetilde P(x) = a_0+a_1x+\cdots+a_nx^n$. Un élément $\alpha\in K$ est
une \textbf{racine} (ou \textbf{zéro}) de $P$ si $\widetilde P(\alpha)=0$.
\end{definition}

\begin{theoreme}[7.2]
Soit $P\in K[X]$ et $\alpha\in K$. Alors :
\begin{enumerate}
  \item Le reste de la division euclidienne de $P$ par $X-\alpha$ est
  $\widetilde P(\alpha)$ (division dont l'existence est établie à la
  section suivante).
  \item $\alpha$ est une racine de $P$ si et seulement si $(X-\alpha)$
  divise $P$.
\end{enumerate}
\end{theoreme}

\begin{proof}
(1) Écrivons $P=(X-\alpha)Q+R$, où $R$ est le reste de la division
euclidienne (Théorème 7.3 ci-après), donc $\deg R < \deg(X-\alpha)=1$,
c'est-à-dire $R$ est un polynôme constant, disons $R=c\in K$. En évaluant
en $\alpha$ : $\widetilde P(\alpha) = (\alpha-\alpha)\widetilde
Q(\alpha)+c = 0+c = c$. Donc $R = \widetilde P(\alpha)$.

(2) Découle immédiatement de (1) : $\alpha$ racine de $P$
$\iff \widetilde P(\alpha)=0 \iff R=0 \iff (X-\alpha)$ divise $P$.
\end{proof}

\begin{definition}[7.4 -- Ordre de multiplicité]
Soit $P\in K[X]\setminus\{0\}$ et $\alpha$ une racine de $P$. L'\textbf{ordre
de multiplicité} de $\alpha$ est le plus grand entier $m\in\mathbb{N}^*$
tel que $(X-\alpha)^m$ divise $P$. Si $m=1$, $\alpha$ est dite
\textbf{racine simple} ; si $m=2$, \textbf{racine double} ; etc.
\end{definition}

\begin{proposition}[7.1 -- Critère de multiplicité par les dérivées]
Soit $P\in K[X]\setminus\{0\}$ et $\alpha\in K$. Alors $\alpha$ est une
racine \textbf{simple} de $P$ si et seulement si
\[
\widetilde P(\alpha) = 0 \qquad \text{et} \qquad \widetilde{P'}(\alpha)
\neq 0,
\]
où $P'$ désigne le polynôme dérivé de $P$.
\end{proposition}

\begin{exemple}[7.1]
Soit $P = X^3+5X^2+3X-9$. Vérifions que $-3$ est racine, et déterminons
son ordre de multiplicité.

$\widetilde P(-3) = -27+45-9-9 = 0$ : $-3$ est bien racine.

$P' = 3X^2+10X+3$, donc $\widetilde{P'}(-3) = 27-30+3=0$ : d'après la
Proposition 7.1, $-3$ \textbf{n'est pas} une racine simple ; c'est (au
moins) une racine double.

On vérifie, par la division euclidienne (ou directement) :
\[
P = (X+3)^2(X-1).
\]
$-3$ est donc une racine \textbf{double}, et $1$ est une racine simple
($\widetilde P(1) = 1+5+3-9=0$).
\end{exemple}

\section{Division euclidienne dans \texorpdfstring{$K[X]$}{K[X]}}


\begin{theoreme}[7.3 -- Division euclidienne dans \texorpdfstring{$K[X]$}{K[X]}]
Soient $P,B\in K[X]$, $B\neq0$. Il existe un \textbf{unique} couple
$(Q,R)\in K[X]^2$ tel que
\[
P = BQ+R \qquad \text{et} \qquad \deg R < \deg B.
\]
\end{theoreme}

\begin{proof}[Démonstration (existence, par récurrence forte sur $\deg P$)]
Notons $m=\deg B$.

\textbf{Si $\deg P < m$ (en particulier si $P=0$) :} on prend $Q=0$,
$R=P$, et $\deg R=\deg P<m=\deg B$ : convient.

\textbf{Si $\deg P = n \geq m$ :} écrivons $P=a_nX^n+\cdots$ et
$B=b_mX^m+\cdots$ (avec $a_n,b_m\neq0$). Posons
\[
P_1 = P - \frac{a_n}{b_m}X^{n-m}B.
\]
Le terme en $X^n$ de $\dfrac{a_n}{b_m}X^{n-m}B$ vaut exactement
$\dfrac{a_n}{b_m}X^{n-m}\times b_mX^m = a_nX^n$ : il s'annule donc avec le
terme dominant de $P$, ce qui donne $\deg P_1 < n$ (ou $P_1=0$).

Par \textbf{hypothèse de récurrence} (le résultat étant supposé vrai pour
tout polynôme de degré $<n$), il existe $Q_1,R\in K[X]$ tels que
$P_1 = BQ_1+R$ avec $\deg R<m$. En substituant :
\[
P = P_1 + \frac{a_n}{b_m}X^{n-m}B = BQ_1+R+\frac{a_n}{b_m}X^{n-m}B =
B\left(Q_1+\frac{a_n}{b_m}X^{n-m}\right) + R,
\]
avec $\deg R<m$ : le couple $Q=Q_1+\dfrac{a_n}{b_m}X^{n-m}$ et $R$
convient.

\textbf{Unicité.} Supposons $P=BQ+R=BQ'+R'$ avec $\deg R<m$ et
$\deg R'<m$. Alors $B(Q-Q')=R'-R$. Si $Q\neq Q'$, en passant aux
degrés : $\deg B + \deg(Q-Q') = \deg(R'-R)$. Or
$\deg(R'-R) \leq \max(\deg R,\deg R') < m = \deg B$, tandis que
$\deg B + \deg(Q-Q') \geq \deg B$ (car $\deg(Q-Q')\geq0$ si
$Q\neq Q'$) : contradiction. Donc $Q=Q'$, puis $R=R'$.
\end{proof}

\begin{exemple}[7.2]
Divisons $A=X^4+X^3-2X+1$ par $B=X^2+1$ :
\[
X^4+X^3-2X+1 = (X^2+1)(X^2+X-1) + (-3X+2).
\]
Vérification du degré du reste : $\deg(-3X+2)=1 < 2 = \deg B$ : correct.
Quotient $Q=X^2+X-1$, reste $R=-3X+2$.
\end{exemple}

\section{Division suivant les puissances croissantes}


\begin{theoreme}[7.4 -- Division suivant les puissances croissantes]
Soient $A,B\in K[X]$ avec $B(0)\neq0$, et $n\in\mathbb{N}$. Il existe un
\textbf{unique} couple $(Q,R)\in K[X]^2$ tel que
\[
A = BQ + X^{n+1}R, \qquad \deg Q \leq n.
\]
$Q$ est appelé le \textbf{quotient à l'ordre $n$} de la division de $A$
par $B$ suivant les puissances croissantes.
\end{theoreme}

\begin{remarque}[Méthode pratique]
On effectue la division comme la division euclidienne usuelle, mais en
ordonnant les polynômes suivant les puissances \textbf{croissantes} de
$X$ (et non décroissantes), et on \textbf{arrête} le processus dès que le
quotient atteint le degré $n$ souhaité.
\end{remarque}

\begin{exemple}[7.3]
Divisons $A=1+2X$ par $B=1-X+X^2$ suivant les puissances croissantes à
l'ordre $n=2$.

On ordonne les deux polynômes suivant les puissances \textbf{croissantes}
de $X$, et on divise comme pour une division euclidienne usuelle, mais en
éliminant à chaque étape le terme de \textbf{plus bas} degré du reste
courant (au lieu du plus haut degré) :

\textbf{Étape 1.} Le terme de plus bas degré de $A$ est $1$, celui de $B$
est $1$ : premier terme du quotient $1$. On calcule $A - 1\times B =
(1+2X) - (1-X+X^2) = 3X-X^2$.

\textbf{Étape 2.} Le terme de plus bas degré de $3X-X^2$ est $3X$ :
deuxième terme du quotient $3X$. On calcule $(3X-X^2) - 3X\times B =
(3X-X^2) - (3X-3X^2+3X^3) = 2X^2-3X^3$.

\textbf{Étape 3.} Le terme de plus bas degré de $2X^2-3X^3$ est $2X^2$ :
troisième terme du quotient $2X^2$. On calcule $(2X^2-3X^3) - 2X^2\times
B = (2X^2-3X^3) - (2X^2-2X^3+2X^4) = -X^3-2X^4$.

Le quotient a atteint le degré $n=2$ souhaité, on arrête : $Q =
1+3X+2X^2$. Le reste de la dernière étape est $-X^3-2X^4 = X^3(-1-2X)$,
d'où, en identifiant avec $X^{n+1}R = X^3R$ : $R = -1-2X$.

\emph{Vérification :}
\[
BQ + X^3R = (1-X+X^2)(1+3X+2X^2) + X^3(-1-2X).
\]
En développant $(1-X+X^2)(1+3X+2X^2) = 1+3X+2X^2-X-3X^2-2X^3+X^2+3X^3+2X^4
= 1+2X+0\cdot X^2+X^3+2X^4$, puis en ajoutant $X^3(-1-2X)=-X^3-2X^4$ :
\[
1+2X+X^3+2X^4-X^3-2X^4 = 1+2X = A.
\]
L'égalité $A=BQ+X^3R$ est bien vérifiée.
\end{exemple}

\section{Factorisation dans \texorpdfstring{$\mathbb{C}[X]$}{C[X]} et \texorpdfstring{$\mathbb{R}[X]$}{R[X]}}

\begin{theoreme}[7.5 -- Théorème de d'Alembert-Gauss (admis)]
Tout polynôme non constant de $\mathbb{C}[X]$ admet au moins une racine
dans $\mathbb{C}$. Par conséquent, tout polynôme de $\mathbb{C}[X]$ de
degré $n\geq1$ se factorise entièrement sous la forme
\[
P = a_n \prod_{k=1}^{r} (X-\alpha_k)^{m_k}, \qquad m_1+\cdots+m_r = n,
\]
où $\alpha_1,\dots,\alpha_r$ sont les racines distinctes de $P$, d'ordres
de multiplicité respectifs $m_1,\dots,m_r$.
\end{theoreme}

\begin{remarque}
La démonstration de ce théorème dépasse largement le cadre de ce cours de
première année (elle utilise des outils d'analyse complexe) ; nous
l'admettons. C'est l'un des résultats les plus importants de toutes les
mathématiques : il affirme que $\mathbb{C}$ est \textbf{algébriquement
clos}.
\end{remarque}

\begin{theoreme}[7.6 -- Factorisation dans \texorpdfstring{$\mathbb{R}[X]$}{R[X]}]
Tout polynôme non constant de $\mathbb{R}[X]$ se factorise en produit de
facteurs de degré $1$ (de la forme $X-\alpha$, $\alpha\in\mathbb{R}$) et
de facteurs de degré $2$ à discriminant strictement négatif (irréductibles
sur $\mathbb{R}$).
\end{theoreme}

\begin{proof}[Idée de démonstration]
Soit $P\in\mathbb{R}[X]\subset\mathbb{C}[X]$. D'après le Théorème 7.5,
$P$ se factorise sur $\mathbb{C}$ en facteurs $(X-\alpha_k)^{m_k}$. Le
point clé est que si $\alpha\in\mathbb{C}\setminus\mathbb{R}$ est une
racine de $P$, alors son conjugué $\bar\alpha$ en est \textbf{aussi} une
racine, de \textbf{même} ordre de multiplicité (car les coefficients de
$P$ sont réels, donc $\widetilde P(\bar\alpha) =
\overline{\widetilde P(\alpha)} = \bar 0 = 0$, en utilisant les propriétés
du conjugué, Chapitre 6). Les racines non réelles de $P$ se regroupent
donc par paires conjuguées $\{\alpha,\bar\alpha\}$, et le produit
\[
(X-\alpha)(X-\bar\alpha) = X^2 - (\alpha+\bar\alpha)X + \alpha\bar\alpha =
X^2 - 2\mathrm{Re}(\alpha)X + |\alpha|^2
\]
est un polynôme à coefficients \textbf{réels}, de discriminant
strictement négatif (car il n'a pas de racine réelle, par construction).
En regroupant ainsi toutes les paires de racines conjuguées, on obtient
la factorisation annoncée.
\end{proof}

\begin{exemple}[7.4]
Factorisons $P=X^4+1$ dans $\mathbb{C}[X]$ puis dans $\mathbb{R}[X]$.

Les racines de $P$ sont les $4$ racines $4$-ièmes de $-1$ (Chapitre 6,
Théorème 6.3) : $-1=e^{i\pi}$, donc les racines sont
$e^{i(\pi+2k\pi)/4}$ pour $k=0,1,2,3$, soit $e^{i\pi/4}$, $e^{i3\pi/4}$,
$e^{i5\pi/4}=e^{-i3\pi/4}$, $e^{i7\pi/4}=e^{-i\pi/4}$. Dans
$\mathbb{C}[X]$ :
\[
X^4+1 = \left(X-e^{i\pi/4}\right)\left(X-e^{-i\pi/4}\right)
\left(X-e^{i3\pi/4}\right)\left(X-e^{-i3\pi/4}\right).
\]
En regroupant les paires conjuguées : $\left(X-e^{i\pi/4}\right)
\left(X-e^{-i\pi/4}\right) = X^2 - 2\cos\!\left(\dfrac{\pi}{4}\right)X+1
= X^2-\sqrt2\,X+1$, et de même
$\left(X-e^{i3\pi/4}\right)\left(X-e^{-i3\pi/4}\right) =
X^2-2\cos\!\left(\dfrac{3\pi}{4}\right)X+1 = X^2+\sqrt2\,X+1$. Donc, dans
$\mathbb{R}[X]$ :
\[
X^4+1 = (X^2-\sqrt2\,X+1)(X^2+\sqrt2\,X+1),
\]
et on vérifie que chacun de ces deux facteurs a un discriminant
strictement négatif ($\Delta = 2-4=-2<0$ dans les deux cas) : ils sont
bien irréductibles sur $\mathbb{R}$.
\end{exemple}

\section{Fractions rationnelles}


\begin{definition}[7.5]
Une \textbf{fraction rationnelle} à coefficients dans $K$ est un quotient
$F=\dfrac{P}{Q}$, où $P,Q\in K[X]$, $Q\neq0$ (deux écritures
$\frac{P}{Q}$ et $\frac{P'}{Q'}$ représentant la même fraction si
$PQ'=P'Q$). L'ensemble des fractions rationnelles à coefficients dans $K$
est noté $K(X)$. Une fraction est dite \textbf{irréductible} si $P$ et
$Q$ n'ont aucune racine commune (forme réduite, obtenue en simplifiant
par les facteurs communs).

Les racines de $Q$ (dans l'écriture irréductible) sont appelées
\textbf{pôles} de $F$ ; l'ordre de multiplicité d'un pôle est celui de la
racine correspondante de $Q$.
\end{definition}

\begin{proposition}[7.2 -- Partie entière]
Soit $F=\dfrac{P}{Q}\in K(X)$. Il existe un \textbf{unique} couple
$(E,R)\in K[X]^2$, avec $\deg R < \deg Q$, tel que
\[
F = E + \frac{R}{Q}.
\]
$E$ est appelée la \textbf{partie entière} (ou partie polynomiale) de
$F$ ; elle est \textbf{non nulle} si et seulement si $\deg P\geq\deg Q$.
\end{proposition}

\begin{proof}
C'est une conséquence immédiate de la division euclidienne
(Théorème 7.3) : on écrit $P=QE+R$ avec $\deg R<\deg Q$, et on divise
par $Q$ : $F=\dfrac{P}{Q} = E+\dfrac{R}{Q}$. L'unicité de $E,R$ découle
de celle de la division euclidienne.
\end{proof}

\section{Décomposition en éléments simples}


\begin{theoreme}[7.7 -- Décomposition en éléments simples sur \texorpdfstring{$\mathbb{C}$}{C} (admis)]
Soit $F=\dfrac{P}{Q}\in\mathbb{C}(X)$ irréductible, avec
$Q = \displaystyle\prod_{i=1}^r (X-\alpha_i)^{m_i}$. Alors $F$ s'écrit de
façon \textbf{unique} :
\[
F = E(X) + \sum_{i=1}^{r} \sum_{k=1}^{m_i} \frac{c_{i,k}}{(X-\alpha_i)^k},
\qquad c_{i,k}\in\mathbb{C},
\]
où $E$ est la partie entière de $F$.
\end{theoreme}

\begin{proposition}[7.3 -- Calcul du coefficient d'un pôle simple]
Si $\alpha$ est un pôle \textbf{simple} de $F=\dfrac{P}{Q}$, le
coefficient $c$ du terme $\dfrac{c}{X-\alpha}$ dans la décomposition est
donné par
\[
c = \lim_{x\to\alpha} (x-\alpha)F(x).
\]
En pratique, si $Q=(X-\alpha)\widetilde Q$ avec $\widetilde
Q(\alpha)\neq0$, on a simplement $c = \dfrac{P(\alpha)}{\widetilde
Q(\alpha)}$.
\end{proposition}

\begin{proof}
Écrivons $F = \dfrac{c}{X-\alpha} + G(X)$, où $G$ regroupe tous les
autres termes de la décomposition (qui n'ont \textbf{pas} $\alpha$ comme
pôle, par unicité de la décomposition). En multipliant par $(X-\alpha)$ :
\[
(X-\alpha)F(X) = c + (X-\alpha)G(X).
\]
Comme $\alpha$ n'est pas un pôle de $G$, $G$ est continue en $\alpha$
(intuitivement -- une justification rigoureuse relève du cours d'analyse),
donc $(X-\alpha)G(X) \to 0$ quand $x\to\alpha$. En passant à la limite :
$\displaystyle\lim_{x\to\alpha}(x-\alpha)F(x) = c$.
\end{proof}

\begin{exemple}[7.5]
Décomposons $F = \dfrac{3X-1}{(X-1)(X+2)}$ en éléments simples sur
$\mathbb{R}$ (les deux pôles $1$ et $-2$ étant réels et simples) :
\[
F = \frac{a}{X-1} + \frac{b}{X+2}.
\]
D'après la Propriété 7.3 :
\[
a = \frac{3\times1-1}{1+2} = \frac{2}{3}, \qquad b =
\frac{3\times(-2)-1}{-2-1} = \frac{-7}{-3} = \frac{7}{3}.
\]
Donc $F = \dfrac{2/3}{X-1} + \dfrac{7/3}{X+2}$.

\emph{Vérification :} $\dfrac{2/3}{X-1}+\dfrac{7/3}{X+2} =
\dfrac{2(X+2)/3+7(X-1)/3}{(X-1)(X+2)} = \dfrac{(2X+4+7X-7)/3}{(X-1)(X+2)}
= \dfrac{(9X-3)/3}{(X-1)(X+2)} = \dfrac{3X-1}{(X-1)(X+2)}$ : correct.
\end{exemple}

\begin{remarque}[Pôle multiple : coefficient du terme de plus haut ordre]
Pour un pôle $\alpha$ d'ordre $m$, le coefficient $c_{i,m}$ du terme
$\dfrac{c_{i,m}}{(X-\alpha)^m}$ (plus haute puissance) se calcule de la
même façon que pour un pôle simple, en multipliant $F$ par $(X-\alpha)^m$
puis en évaluant en $\alpha$. Les coefficients des puissances
inférieures s'obtiennent ensuite en évaluant l'égalité (après
multiplication par $(X-\alpha)^m$) en d'autres valeurs bien choisies, ou
par identification.
\end{remarque}

\begin{exemple}[7.6 -- Pôle double]
Décomposons $F = \dfrac{1}{(X-1)^2(X+1)}$ en éléments simples :
\[
F = \frac{a}{X-1} + \frac{b}{(X-1)^2} + \frac{c}{X+1}.
\]
\textbf{Coefficient $b$ (pôle double, puissance maximale) :} en
multipliant par $(X-1)^2$ et en évaluant en $X=1$ :
\[
b = \left.\frac{1}{X+1}\right|_{X=1} = \frac12.
\]
\textbf{Coefficient $c$ (pôle simple) :} d'après la Propriété 7.3,
\[
c = \left.\frac{1}{(X-1)^2}\right|_{X=-1} = \frac{1}{4}.
\]
\textbf{Coefficient $a$ :} on évalue l'égalité en une valeur commode,
par exemple $X=0$ :
\[
F(0) = \frac{1}{(0-1)^2(0+1)} = 1, \qquad \text{et} \qquad
\frac{a}{-1}+\frac{b}{1}+\frac{c}{1} = -a+\frac12+\frac14 = -a+\frac34.
\]
En égalant : $1 = -a+\dfrac34$, donc $a = \dfrac34-1 = -\dfrac14$.

Ainsi $F = \dfrac{-1/4}{X-1} + \dfrac{1/2}{(X-1)^2} + \dfrac{1/4}{X+1}$.
\end{exemple}

\begin{theoreme}[7.8 -- Décomposition en éléments simples sur \texorpdfstring{$\mathbb{R}$}{R} (admis)]
Soit $F=\dfrac{P}{Q}\in\mathbb{R}(X)$ irréductible, avec $Q$ factorisé
dans $\mathbb{R}[X]$ (Théorème 7.6) en facteurs $(X-\alpha_i)^{m_i}$ et
$(X^2+p_jX+q_j)^{n_j}$ (facteurs irréductibles). Alors $F$ s'écrit de
façon unique comme somme de sa partie entière, de termes
$\dfrac{c}{(X-\alpha_i)^k}$ ($c\in\mathbb{R}$), et de termes
$\dfrac{aX+b}{(X^2+p_jX+q_j)^k}$ ($a,b\in\mathbb{R}$).
\end{theoreme}

\begin{remarque}[Méthode pratique sur $\mathbb{R}$]
Une méthode efficace consiste à décomposer $F$ sur $\mathbb{C}$
(Théorème 7.7), où tous les pôles sont linéaires, puis à \textbf{regrouper}
les termes associés à des pôles complexes conjugués : la somme de deux
tels termes redonne automatiquement un terme réel de la forme
$\dfrac{aX+b}{X^2+pX+q}$.
\end{remarque}

\begin{exemple}[7.7]
Décomposons $F=\dfrac{1}{X^2+1}$ sur $\mathbb{C}$ puis sur $\mathbb{R}$.

\textbf{Sur $\mathbb{C}$ :} $X^2+1=(X-i)(X+i)$, pôles simples $i,-i$.
D'après la Propriété 7.3 :
\[
c_i = \left.\frac{1}{X+i}\right|_{X=i} = \frac{1}{2i} = -\frac{i}{2},
\qquad c_{-i} = \left.\frac{1}{X-i}\right|_{X=-i} = \frac{1}{-2i} =
\frac{i}{2}.
\]
Donc $F = \dfrac{-i/2}{X-i} + \dfrac{i/2}{X+i}$ dans $\mathbb{C}(X)$.

\textbf{Sur $\mathbb{R}$ :} le polynôme $X^2+1$ est déjà irréductible sur
$\mathbb{R}$ (discriminant $-4<0$), donc $F$ est \textbf{déjà} un élément
simple sur $\mathbb{R}$ : sa décomposition réelle est $F$ elle-même,
$F = \dfrac{0\times X+1}{X^2+1}$. Ceci illustre que la décomposition sur
$\mathbb{R}$ peut être \og plus grossière \fg{} que sur $\mathbb{C}$,
précisément à cause des pôles non réels qui se regroupent.
\end{exemple}

\begin{exemple}[7.8]
Décomposons $F=\dfrac{1}{X(X^2+1)}$ sur $\mathbb{R}$ :
\[
F = \frac{a}{X} + \frac{bX+c}{X^2+1}.
\]
\textbf{Coefficient $a$ (pôle réel simple) :} d'après la Propriété 7.3,
$a = \left.\dfrac{1}{X^2+1}\right|_{X=0} = 1$.

\textbf{Coefficients $b,c$ :} on soustrait la partie déjà connue :
\[
F - \frac{1}{X} = \frac{1-(X^2+1)}{X(X^2+1)} = \frac{-X^2}{X(X^2+1)} =
\frac{-X}{X^2+1}.
\]
Donc $\dfrac{bX+c}{X^2+1} = \dfrac{-X}{X^2+1}$, d'où $b=-1$, $c=0$.

Ainsi $F = \dfrac1X - \dfrac{X}{X^2+1}$.
\end{exemple}

\section{Applications}

\begin{application}[Calcul de primitives]
La décomposition en éléments simples est la méthode \textbf{universelle}
pour calculer la primitive d'une fraction rationnelle : chaque terme
$\dfrac{c}{(X-\alpha)^k}$ se primitive immédiatement (logarithme si
$k=1$, puissance sinon), et chaque terme quadratique se ramène à un
arctangente après complétion du carré.
\end{application}

\begin{application}[Automatique et traitement du signal]
En automatique, la \textbf{fonction de transfert} d'un système linéaire
est une fraction rationnelle en la variable de Laplace $p$ ; sa
décomposition en éléments simples permet de calculer explicitement la
réponse temporelle du système (chaque pôle correspond à un mode
exponentiel ou oscillant).
\end{application}

\begin{application}[Codes correcteurs et calcul formel]
L'arithmétique de $K[X]$ (division euclidienne, factorisation) est à la
base des \textbf{codes de Reed-Solomon} (utilisés dans les CD, DVD,
codes-barres QR) et de tous les logiciels de \textbf{calcul formel}
(manipulation exacte de polynômes et fractions rationnelles).
\end{application}

\section{Résumé du chapitre}

\begin{resume}
\begin{itemize}
  \item $(K[X],+,\cdot)$ est un anneau commutatif intègre ;
  $\deg(PQ)=\deg P+\deg Q$ (Théorème 7.1).
  \item $\alpha$ racine de $P$ $\iff$ $(X-\alpha)\mid P$ (Théorème 7.2) ;
  ordre de multiplicité via $(X-\alpha)^m \mid P$, ou via les dérivées
  successives (Propriété 7.1).
  \item \textbf{Division euclidienne} dans $K[X]$ : $P=BQ+R$,
  $\deg R<\deg B$ (Théorème 7.3), analogue exact du Chapitre 5.
  \item \textbf{Division suivant les puissances croissantes} : $A=BQ+X^{n+1}R$,
  $\deg Q\leq n$ (Théorème 7.4), utile pour les développements au
  voisinage de $0$.
  \item \textbf{d'Alembert-Gauss} (admis) : $\mathbb{C}$ est
  algébriquement clos ; factorisation complète sur $\mathbb{C}$, en
  facteurs de degré $1$ ou $2$ (discriminant $<0$) sur $\mathbb{R}$.
  \item \textbf{Décomposition en éléments simples} : toute fraction se
  décompose (de façon unique) en partie entière + termes
  $\dfrac{c}{(X-\alpha)^k}$ (sur $\mathbb{C}$, ou sur $\mathbb{R}$ pour
  les pôles réels) + termes $\dfrac{aX+b}{(X^2+pX+q)^k}$ (sur
  $\mathbb{R}$, pour les paires de pôles complexes conjugués). Coefficient
  d'un pôle simple $\alpha$ : $c=\dfrac{P(\alpha)}{\widetilde Q(\alpha)}$
  où $Q=(X-\alpha)\widetilde Q$.
\end{itemize}
\end{resume}

\end{document}
